\documentclass[11pt,a4paper]{article}
\usepackage[T1]{fontenc}
\usepackage[utf8]{inputenc}
\usepackage{lmodern}
\usepackage[margin=25mm,headheight=15pt]{geometry}
\usepackage{amsmath,amssymb,amsthm,mathtools}
\usepackage{microtype}
\usepackage{booktabs,longtable,array}
\usepackage{xcolor}
\usepackage{xurl}
\usepackage{enumitem}
\usepackage{fancyhdr}
\usepackage{hyperref}
\definecolor{ink}{HTML}{17324D}
\definecolor{accent}{HTML}{156A72}
\hypersetup{colorlinks=true,linkcolor=ink,citecolor=accent,urlcolor=accent,
 pdftitle={Category Theory in Intuitionistic Primitive Recursive Arithmetic},
 pdfauthor={Category theory research report},
 pdfsubject={Constructive foundations, explicit proofs, and an object-language Lean roadmap}}
\newcommand{\ipra}{\ensuremath{\mathsf{iPRA}}}
\newcommand{\N}{\mathbb{N}}
\newcommand{\Hom}{\operatorname{Hom}}
\newcommand{\Nat}{\operatorname{Nat}}
\newcommand{\id}{\operatorname{id}}
\newcommand{\op}{\mathrm{op}}
\newcommand{\Set}{\mathsf{Set}}
\newcommand{\Ob}{\operatorname{Ob}}
\newcommand{\Arr}{\operatorname{Arr}}
\theoremstyle{plain}
\newtheorem{theorem}{Theorem}[section]
\newtheorem{proposition}[theorem]{Proposition}
\newtheorem{lemma}[theorem]{Lemma}
\newtheorem{corollary}[theorem]{Corollary}
\theoremstyle{definition}
\newtheorem{definition}[theorem]{Definition}
\newtheorem{example}[theorem]{Example}
\theoremstyle{remark}
\newtheorem{remark}[theorem]{Remark}
\numberwithin{equation}{section}
\setlist{topsep=5pt,itemsep=3pt,leftmargin=*}
\setlength{\parskip}{3pt}
\setlength{\emergencystretch}{2em}
\setcounter{tocdepth}{2}
\pagestyle{fancy}
\fancyhf{}
\fancyhead[L]{\small\color{ink}Category theory in \ipra}
\fancyhead[R]{\small 20 November 2026}
\fancyfoot[C]{\thepage}
\renewcommand{\headrulewidth}{0.3pt}
\begin{document}
\hypersetup{pageanchor=false}
\begin{titlepage}
\thispagestyle{empty}
\setlength{\parindent}{0pt}
\vspace*{17mm}
{\Large\color{accent}\sffamily CONSTRUCTIVE FOUNDATIONS\par}
\vspace{9mm}
{\Huge\bfseries\color{ink}Category Theory in\\[3mm]
Intuitionistic Primitive\\[3mm]Recursive Arithmetic\par}
\vspace{9mm}
{\Large Explicit mathematical proofs and\\[2mm]
an object-language Lean formalization path\par}
\vspace{12mm}
{\large Research report and reproducibility package\par}
\vspace{2mm}
{\large 20 November 2026\par}
\vfill
\noindent\rule{\linewidth}{0.6pt}
\begin{abstract}
A substantial structural core of category theory can be developed in
first-order intuitionistic primitive recursive arithmetic by fixing
arithmetic object and hom predicates, an arithmetic equivalence on
parallel arrows, and primitive recursive code operations. The proofs
may use undecidable equality formulas while remaining finite
intuitionistic derivations. We give full human-readable proofs of
Yoneda, representability and density; selected finite limits and colimits;
adjunctions, preservation and split equivalences; idempotent completion;
monads, Kleisli and Eilenberg--Moore constructions; pointwise Kan
extensions under explicit hypotheses; and a constructive diagonal
theorem. Uniform finite-table algorithms separately provide all finite
limits, colimits and exponentials, including empty sets.
The report distinguishes iPRA statements from external schemas,
selected structure from bare existence, and reused upstream
proof source from new uncompiled Lean integration drafts.
\end{abstract}
\noindent\rule{\linewidth}{0.6pt}\par
\vspace{3mm}
{\small\color{ink}
Package includes the \LaTeX{} source, PDF, exact finite-model code,
Lean drafts, a pinned upstream import closure, proof-status records,
and reproduction instructions.\par}
\end{titlepage}
\hypersetup{pageanchor=true}
\tableofcontents
\clearpage
\section{Main results and how to read them}\label{sec:overview}

The answer is positive for a broad, precisely specified core of category
theory. Its structural arguments use very little arithmetic: the real work
is to give their data a legitimate arithmetic interface. One can permit
arithmetic, potentially undecidable equality of arrows without allowing
arbitrary function spaces or induction on all arithmetic formulas.

This report uses the user's overview \cite{wiki} to delimit the basic
subject, and standard author-hosted references for terminology
\cite{riehl,awodey}. Its mathematical arguments are written out in full.
The classical theorems themselves are not claimed to be new. The purpose
is a coherent iPRA realization with every computational and logical
boundary visible, building on the pinned ProveIt development.

\subsection{The principal theorem families}
\begin{longtable}{@{}>{\raggedright\arraybackslash}p{0.27\linewidth}>{\raggedright\arraybackslash}p{0.60\linewidth}p{0.06\linewidth}@{}}
\toprule
Result & Exact form established here & Sec.\\
\midrule
\endhead
Arithmetic foundations & Fixed formulas for objects, arrows and equality;
PR code operations; specialization of finite structural proofs. &
\ref{sec:foundation}\\
Functorial calculus & Functors, natural transformations, whiskering,
horizontal and vertical composition, interchange, duality. &
\ref{sec:functors-yoneda-foundation}\\
Yoneda & Explicit mutually inverse constructions
$\alpha\mapsto\alpha_A(\id_A)$ and $x\mapsto(f\mapsto P(f)x)$,
with all typing and naturality laws. & \ref{sec:yoneda}\\
Representability & Full faithfulness, canonical uniqueness of representing
objects, and representation by a chosen universal element. &
\ref{sec:yoneda}\\
Density / co-Yoneda & Canonical generator normal form and a universal
cocone, with arithmetic equality defined by equality of values. &
\ref{sec:density}\\
Universal constructions & Unique isomorphisms; pullbacks; all fixed finite
limits from selected terminal objects, products and equalizers, plus duals. &
\ref{sec:limits}\\
Finite-set existence & Uniform PR algorithms for finite limits, colimits
and exponentials, including $[0]$ and $0^0=1$. & \ref{sec:finite}\\
Adjunctions & Hom correspondences iff unit/counit triangles; uniqueness;
preservation of supplied limits and colimits. & \ref{sec:adjunctions}\\
Equivalence & A quasi-inverse constructed from faithful, split full and
split essentially surjective data. & \ref{sec:adjunctions}\\
Monads & Monad of an adjunction; Kleisli and Eilenberg--Moore categories
and adjunctions; creation of supplied algebra limits. & \ref{sec:monads}\\
Idempotents & Canonical splitting in the Karoubi envelope and extension
along displayed splittings. & \ref{sec:limits}\\
Kan extensions & Pointwise formula with supplied uniform comma
(co)limits; unconditional existence for finite tables in finite sets. &
\ref{sec:kan}\\
Diagonal theorem & Lawvere fixed points and the obstruction to an
unrestricted PR evaluation family. & \ref{sec:lawvere}\\
\bottomrule
\end{longtable}

\subsection{Three levels of assertion}
A result about a fixed presentation is an ordinary arithmetic theorem.
A result quantified over a numerical family of finite tables is also an
ordinary arithmetic theorem, now with table codes as variables.
A result ranging over arbitrary presentation syntax is an external
schema producing such theorems. These levels are distinguished throughout.

The structural theorems are relative to the explicitly stated category
laws and selected universal operations. When those laws have iPRA proofs,
the theorem's proof specializes to an iPRA proof. This is mathematically
substantial: the Yoneda and adjunction calculations apply also to infinite
presented categories with undecidable arrow equality. They are not
restricted to testing finite tables.

Conversely, an existence claim that a category has all limits
requires more than a formal calculation. The finite-set section supplies
actual constructors and a quantifier-free induction analysis. For
arbitrary presentations, the report states the exact universal operations
needed and does not invent them.

\subsection{What the deliverable's verification labels mean}
\begin{description}[style=nextline]
\item[Human proof]
The displayed assumptions and conclusions are connected by complete
mathematical arguments in this report. Numerical algorithms include
their invariants and primitive-recursive bounds; structural proofs use
only finite intuitionistic inferences.
\item[Reused object-language source]
The pinned ProveIt modules actually construct judgments of the form
$\Gamma\vdash[\mathsf{PRAAx}]\varphi$. They were read and their import
closure included. The repository reports its derivation library as
kernel-checked; it was not rebuilt in this environment.
\item[New Lean integration draft]
The three new modules target that same derivation syntax and contain
no proof placeholders, but were not compiled because the Lean executable
was unavailable. Their elaboration and kernel checking remain to be checked.
\item[Exact finite-model check]
A reproducible Python run checks the stated finite cases with exact
integer arithmetic. It supports implementation validation but is not
an iPRA formalization.
\end{description}

Useful refinements relative to the inspected upstream sources include
an interface for Yoneda at one fixed parameter and representing object;
a direct density presentation that avoids arbitrary quotient existence;
a complete empty-set treatment and explicit uniform finite-diagram
tester; and one coherent proof and formalization plan for the monad,
algebra, Kan-extension and idempotent constructions.
No claim of a completed Lean mechanization of all these refinements is made.


\section{The arithmetic foundation and its exact meaning}\label{sec:foundation}

\subsection{The base theory}
Throughout, \ipra\ means first-order intuitionistic primitive recursive arithmetic.
Its only variables range over natural numbers. Its terms contain $0$, successor,
and a symbol for each fixed primitive recursive (PR) definition; symbols come
with their defining equations. Its nonlogical axioms are these equations,
successor injectivity, $Sx\ne0$, and the universal closures of
\[
 [A(0,\vec u)\wedge\forall n(A(n,\vec u)\to A(Sn,\vec u))]
       \longrightarrow\forall n A(n,\vec u),                 \tag{I}
\]
where $A$ is quantifier-free. Logical rules are the usual intuitionistic
introduction and elimination rules for $\bot,\wedge,\vee,\to,\forall,\exists$,
with equality reflexivity and substitution. Negation abbreviates implication
to $\bot$. Quantifiers in a theorem need not be bounded or quantifier-free;
only the induction formula in (I) is restricted. This is also the base theory
specified by ProveIt's current iPRA sources \cite{repo-readme,repo-pr}.

We use PR encodings of tuples, finite lists, and finite tables. For example,
pairing and its projections may be defined by bounded search within an
elementary numerical bound, and a list may be coded by its length and a
base large enough for its entries. All searches in the finite algorithms below
come with a specified numerical bound. Default values make these operations
total on malformed inputs; no categorical conclusion concerns those defaults.
Arithmetic order, Boolean tests, finite sums and products, finite-list lookup,
and bounded minimization have their familiar PR definitions.

\begin{lemma}[Bounded arithmetic bookkeeping]\label{lem:bounded}
A fixed formula using PR predicates, Boolean connectives and bounded numerical
quantifiers has a PR characteristic function. Bounded scans, table filtering,
and finite-list updates with PR length bounds are PR. The loop invariants used
for these constructions can be expressed quantifier-free and proved by their
indicated base and step calculations.
\end{lemma}
\begin{proof}
Use $1$ for true and $0$ for false. Conjunction is multiplication, negation is
$1-b$, and bounded universal quantification is an iterated product of bits;
bounded existential quantification is one minus the product of complementary
bits. All are primitive recursions. A bounded least witness is obtained by a
scan carrying a flag recording whether an earlier witness was found and its
index. A table update is a scan rebuilding its coded entries.
The defining equations prove the base and step equations for these scans.
For example, the assertion that a returned index is the first successful
one is expressed by its range, its test bit, and a product of failure bits
at earlier indices. This assertion is quantifier-free after replacing the
bounded product by its PR symbol. The same device expresses every bounded
table invariant used below. Consequently induction is applied to a numerical
equation or a Boolean combination of such equations, never to an assertion
about all arbitrary functions. The correspondence between bounded quantifier
notation and its Boolean test follows from the same scan equations, with
any currently discussed index treated as an additional numerical parameter.
\end{proof}

\subsection{Categories as fixed arithmetic presentations}
\begin{definition}[Presented category]\label{def:category}
A presentation $C$ consists of fixed arithmetic formulas
\[
 O_C(a),\qquad H_C(a,b,f),\qquad E_C(a,b,f,g)
\]
and fixed PR terms $\id_C(a)$ and $\operatorname{cmp}_C(a,b,c,g,f)$.
The formula $H_C(a,b,f)$ means that $f$ is an arrow from $a$ to $b$.
The formula $E_C(a,b,f,g)$ compares parallel arrow codes, and is not
necessarily numerical equality or a decidable relation. We write $f\sim g$
for this relation, or simply $f=g$ in categorical equations with the hom
type understood. Write $gf$ or $g\circ f$ for $\operatorname{cmp}(a,b,c,g,f)$.

The presentation comes with iPRA proofs, or explicitly retained hypotheses,
of the following laws. Hom membership entails object membership. On every
hom domain $E_C$ is reflexive, symmetric and transitive. Identity and
composition are well typed. Composition respects $E_C$ in both arguments.
For composable arrows,
\[
 h(gf)\sim(hg)f,\qquad \id_b f\sim f,\qquad f\id_a\sim f.       \tag{C}
\]
All these laws are universally closed implications with the indicated
typing premises. We also require equality judgments to concern well-typed
parallel arrows, or include those premises explicitly whenever they are used.
\end{definition}

The definitions are abbreviations in arithmetic, not new set-existence
axioms. Objects have literal numerical codes; isomorphic objects need not
have equal codes. Only arrows are identified by their hom equivalence.
Different codes for the same informal structure may therefore coexist.
An empty category is allowed, since $O_C$ may be false everywhere.

There are three distinct uses of this definition.
\begin{enumerate}
\item A \emph{fixed presentation}: its formulas and PR terms are chosen
before an arithmetic sentence is formed. Theorems about it are ordinary
iPRA sentences.
\item A \emph{uniform numerical family}: the same formulas and terms have
an extra parameter $p$. The family can be quantified over $p$ because
evaluation of those fixed terms is already PR. Finite table categories
are a particularly useful family of this kind.
\item An \emph{external collection of presentations}: Lean or the
mathematician ranges over formula syntax and PR term syntax. A theorem
constructs an iPRA proof for each instance. This is a proof schema, not
an internal universe whose elements evaluate all PR programs.
\end{enumerate}
Whenever the report says ``every presented category'', the third meaning
is intended unless a numerical family is explicitly supplied.

\begin{theorem}[Specialization of structural proofs]\label{thm:export}
Fix presentations, functors, families and selected universal operations as
defined in this report. A finite typed proof obtained by intuitionistic
logical rules, the supplied laws, and substitution of their PR terms becomes
an iPRA derivation after expanding the abbreviations. If its laws have iPRA
proofs, the conclusion has an iPRA proof; if they are hypotheses, the
conclusion is a relative iPRA theorem with those hypotheses.
No additional induction principle is needed for such a proof.
\end{theorem}
\begin{proof}
Replace every object variable by a number variable with guard $O_C$ and
every arrow variable by a number variable with its hom guard. Replace
categorical equality by $E_C$, and categorical operations by their fixed
terms. A well-typed use of reflexivity, symmetry, transitivity, congruence,
associativity or identity is exactly an instance of a supplied law.
Composition of proof steps is implication elimination or cut. An explicit
witness term justifies existential introduction; an existential hypothesis
may be eliminated locally with the usual eigenvariable condition.
Generalization introduces only numerical quantifiers and their guards.
Thus each line translates to an allowed inference from earlier lines.
There are finitely many lines in a fixed proof, so this construction needs
no induction in the object theory. Formally, induction on proof syntax is
performed in the metatheory to describe the translator, not admitted as
an extra object-language rule. Every induction in a concrete arithmetic
implementation must separately satisfy (I).
\end{proof}

This theorem is intentionally a specialization theorem. It does not assert
that arbitrary Lean theorems are iPRA theorems. In particular, a Lean field
$\N\to\N$ does not certify a PR implementation, and a proof using induction
on a proposition containing unbounded quantifiers does not pass this audit
merely because its conclusion is arithmetical. ProveIt's
\texttt{SetoidInterpretation.sound} gives a more general syntactic
interpretation mechanism when all domain and equality compatibility laws
are supplied \cite{repo-setoid}. The present proofs use the smaller explicit
typed calculus just described.

\subsection{Examples that do not reduce the subject to finite categories}
Every fixed PR monoid $(M,\cdot,e)$ gives a one-object presentation: arrows
are its coded elements, identity is $e$, composition is multiplication.
The additive monoid of natural numbers is the elementary example already
instantiated in \texttt{PresentedCategory.natAdd\_laws}.
A PR preorder gives a thin category: take $H(a,b,f)$ to assert the order
relation and $f=0$; all parallel arrows are equal.
Integer matrices give an infinite presentation with dimension codes as
objects, finite coded matrices as arrows and matrix multiplication as
composition. Its associativity reduces to finite distributivity and
reindexing of bounded sums, hence to Lemma~\ref{lem:bounded} and the PR
ring laws. This example has infinitely many objects and arrows.

Equality may also have unbounded logical complexity. For instance,
the equality of two represented numerical sequences may be
$\forall n\,a(n)=b(n)$, and regular real names may be compared by
an arithmetical approximation condition. For a fixed family of PR names,
operations and a modulus, these are legitimate formulas. Establishing
the analytic category laws is a separate mathematical task; the structural
theorems do not decide equality or assert those laws without proof.

\subsection{Relational programs: a useful extension with a different interface}
There is a further way to have numerical program codes without a total
PR evaluator. Fix a deterministic register-machine language and a PR
predicate $T(e,x,y,t)$ saying that $t$ is a complete halting computation
of code $e$ on input $x$ with output $y$. Its finite-transition check is
PR. Let
\[
 \operatorname{Tot}(e)\equiv\forall x\,\exists y,t\,T(e,x,y,t).
\]
Take one object, arrows $e$ satisfying $\operatorname{Tot}(e)$, and
\[
 e\sim e'\quad\Longleftrightarrow\quad
 \forall x,y,z,t,u\,[T(e,x,y,t)\wedge T(e',x,z,u)\to y=z].       \tag{R}
\]
Identity and composition are PR \emph{code constructors}. To make the
presentation fully explicit, composition first runs its first component,
then runs its second component on the result, with disjoint renamed
work registers. Its trace certificate stores the two component traces
and their joining configuration. Inlining these tagged certificates into
ordinary transition traces is a PR operation.

\begin{proposition}[Relational composition]
Given the finite trace-composition and decomposition lemmas for this fixed
machine language, (R) defines a presented category. The proof uses neither
choice nor an internal assertion that all program codes terminate.
\end{proposition}
\begin{proof}
Determinism gives reflexivity of (R), and symmetry is immediate. For
transitivity on total arrows, obtain an output trace of the middle program
at the particular input $x$; the two equality hypotheses compare the first
and last outputs with that middle output. This is one existential
elimination, not a simultaneous selection over all inputs.
For totality of a composite, fix $x$, use totality of $f$ to obtain $y,t$,
then totality of $g$ at that $y$ to obtain $z,u$, and concatenate the two
traces. Equality is preserved by composition by applying (R) first to
the input-side results and then to the output-side results. Both
associations of three programs run the same three computations; decomposing
their traces and using determinism proves equality. Identity traces give
the identity laws. Each argument compares only finitely many runs at a
time. The trace lemmas themselves are proved by quantifier-free induction
on coded finite traces and their bounded local checks.
\end{proof}

This optional interface is a precise conditional construction: the report
does not ship a newly verified register-machine trace library. It explains
why the evaluator obstruction below does not rule out arithmetically
defined collections of total programs. It also explains the limitation:
$\operatorname{Tot}(e)$ is not a PR test, and extracting a table of outputs
for a variable number of such runs requires its own trace, clock, or
batch-totality proof. The structural development uses fixed PR operations
on arrow codes, not a claim that semantic application is itself PR.

\subsection{Elementary consequences and categorical constructions}
An isomorphism $f:a\to b$ consists of a specified $g:b\to a$ with
$gf=\id_a$ and $fg=\id_b$. A monomorphism cancels on the left and an
epimorphism cancels on the right; these are guarded arithmetic formulas.

\begin{proposition}\label{prop:elementary}
In every presented category inverses are unique. Every section is monic,
every retraction is epic, and every monic retraction is an isomorphism.
Functors preserve isomorphisms; faithful functors reflect equality of
parallel arrows.
\end{proposition}
\begin{proof}
If $gf=\id_a$ and $fh=\id_b$, then
$g=g\id_b=g(fh)=(gf)h=\id_a h=h$.
If $rs=\id_a$ and $su=sv$, then
$u=(rs)u=r(su)=r(sv)=(rs)v=v$.
Reversing arrows gives the retraction assertion.
If $fr=\id_b$ and $f$ is monic, then
$f(rf)=(fr)f=f=f\id_a$, so $rf=\id_a$.
Applying a functor to the two inverse equations gives the inverse equations
for its image. Reflection of equality is exactly faithfulness.
\end{proof}
Being both monic and epic does not in general imply invertibility: in the
thin category on $0<1$, the arrow $0\to1$ cancels on both sides, but there
is no reverse arrow. This finite counterexample is already available in iPRA.

The opposite category reverses hom domains and composition, leaving all
codes unchanged. Its laws are the reversed laws of $C$. Reversing a fixed
proof line by line therefore proves the dual theorem. The product of two
categories pairs object and arrow codes and acts componentwise.
For a fixed object $b$, the slice $C/b$ has objects $(a,u:a\to b)$ and
arrows $f:(a,u)\to(a',u')$ satisfying $u'f=u$.
The coslice is dual. The comma category $(F\downarrow G)$ has objects
$(a,b,u:Fa\to Gb)$ and arrows $(f:a\to a',g:b\to b')$ satisfying
$Gg\,u=u'\,Ff$. Equality is componentwise.
For identity the square is an identity equation. For composition, if
$Gg\,u=u'Ff$ and $Gg'\,u'=u''Ff'$, then
$G(g'g)u=Gg'(Gg\,u)=Gg'(u'Ff)=u''F(f'f)$.
Associativity and both unit laws are componentwise laws. Thus all these
constructions have PR code operations and arithmetic domain formulas;
no collection of arbitrary sets or functions was formed.

\section{Functors, naturality, and the meaning of schematic categories}
\label{sec:functors-yoneda-foundation}

Throughout this section, $C$ and $D$ are fixed arithmetic presentations of
categories.  Objects and arrows have natural-number codes.  Each presentation
has arithmetic formulas for validity and arrow equality, and fixed primitive
recursive operations for identities and composition.  We write $f:a\to b$
for the full validity condition and $f\sim g$ for its hom-setoid equality.
The category laws include preservation of typing and of $\sim$ by every
operation.  In particular, $\sim$ need not be numerical equality or a decidable
relation.  All equations below mean this specified equality, and all variables
are accompanied by their displayed typing hypotheses.  Operations may take
the displayed endpoint codes as additional arguments; those arguments are
suppressed in categorical notation.

\subsection{Functors and transformations}

\begin{definition}[Presented functor]
A functor $F:C\to D$ consists of fixed primitive recursive functions $F_0$
and $F_1$ on object and arrow codes.  Its laws say that $F_0$ preserves object
validity, that $f:a\to b$ implies $F_1(f):F_0(a)\to F_0(b)$, that $F_1$
respects arrow equality, and that
\[
 F_1(\id_a)\sim\id_{F_0(a)},\qquad
 F_1(g\circ f)\sim F_1(g)\circ F_1(f).
\]
We suppress the subscripts.  The laws are arithmetic formulas, proved in
$\ipra$ for a concrete presentation, or explicit premises of a relative
structural theorem.
\end{definition}

\begin{definition}[Presented natural transformation]
For fixed functors $F,G:C\to D$, a natural transformation $\alpha:F\Rightarrow G$
is a fixed primitive recursive component function $a\mapsto\alpha_a$ such that
$\alpha_a:Fa\to Ga$ and, for every $f:a\to b$,
\begin{equation}
 G(f)\circ\alpha_a\sim\alpha_b\circ F(f).
 \label{eq:ordinary-naturality}
\end{equation}
Two transformations are equal when their components are equal at every valid
object.  Thus transformation equality is itself an arithmetic universal
formula; it is not a decision procedure for functions.
\end{definition}

\begin{proposition}[The elementary naturality calculus]
Identity functors, composites of functors, identity transformations, vertical
composition, and whiskering are given by primitive recursive operations on
their displayed data.  They satisfy the category and interchange laws
schematically.
\end{proposition}
\begin{proof}
The identity functor leaves each code unchanged.  A composite functor applies
its two fixed object and arrow functions successively.  Preservation of
identity and composition follows by applying the outer functor's congruence
law to the inner functor's corresponding law, and then applying the outer
law.  Associativity of these composites is equality of the displayed terms.

The identity transformation on $F$ has component $\id_{Fa}$.  Its naturality
square commutes by the two identity laws.  If $\alpha:F\Rightarrow G$ and
$\beta:G\Rightarrow H$, define
$(\beta\alpha)_a=\beta_a\circ\alpha_a$.  Its naturality is the calculation
\[
\begin{split}
 H(f)\circ(\beta_a\circ\alpha_a)
 &\sim (H(f)\circ\beta_a)\circ\alpha_a\\
 &\sim (\beta_b\circ G(f))\circ\alpha_a\\
 &\sim\beta_b\circ(G(f)\circ\alpha_a)\\
 &\sim\beta_b\circ(\alpha_b\circ F(f))\\
 &\sim(\beta_b\circ\alpha_b)\circ F(f).
\end{split}
\]
The identity and associativity laws for vertical composition hold component
by component.  Postwhiskering by $K:D\to E$ has components $K(\alpha_a)$;
its laws follow from functoriality.  Prewhiskering by $L:B\to C$ has components
$\alpha_{L b}$; its laws are instances of naturality for arrows $L(f)$.

For clarity, horizontal composition is also explicit.  If
$\alpha:F\Rightarrow G:C\to D$ and $\beta:H\Rightarrow K:D\to E$, put
\[
 (\beta*\alpha)_a=\beta_{Ga}\circ H(\alpha_a).
\]
Naturality of $\beta$ at $\alpha_a$ identifies this with
$K(\alpha_a)\circ\beta_{Fa}$.  Suppose also
$\alpha':G\Rightarrow J$ and $\beta':K\Rightarrow L$.  The component of the
vertical composite of the two horizontal composites is
\[
 \beta'_{Ja}\circ K(\alpha'_a)\circ\beta_{Ga}\circ H(\alpha_a).
\]
Naturality of $\beta$ at $\alpha'_a$ changes the middle two factors to
$\beta_{Ja}\circ H(\alpha'_a)$.  Functoriality of $H$ then gives
\[
 (\beta'_{Ja}\circ\beta_{Ja})\circ H(\alpha'_a\circ\alpha_a),
\]
which is the component of $(\beta'\beta)*(\alpha'\alpha)$.  Parentheses may
be changed by associativity.  This proves interchange.  Every argument just
given is a finite equational derivation, with no use of induction or excluded
middle.
\end{proof}

This result does \emph{not} assert the existence inside arithmetic of a
category whose arrows are all primitive recursive functors or all primitive
recursive transformations.  It supplies a theorem for each fixed finite
configuration of such data.  A separately specified, uniformly evaluable
coded family can itself be made into a presented category if the required
closure and evaluation operations are supplied.  The word ``schematic''
will retain this meaning below.

\subsection{Setoid-valued presheaves without a universe of sets}

\begin{definition}[Presented presheaf]
A presheaf $P$ on $C$ consists of the following data:
\begin{enumerate}
\item an arithmetic formula $P_0(a,x)$ specifying valid elements in its fibre
      at $a$;
\item an arithmetic equivalence relation $x\approx_a y$ on that fibre;
\item a fixed primitive recursive action function, written $P(f)x$, which
      takes $f:a\to b$ and $x\in P(b)$ to an element of $P(a)$.
\end{enumerate}
The action respects both arrow equality and fibre equality, and satisfies
\begin{equation}
 P(\id_a)x\approx_a x,\qquad
 P(g\circ f)x\approx_a P(f)(P(g)x)
 \quad(f:a\to b, g:b\to c, x\in P(c)).
 \label{eq:presheaf-laws}
\end{equation}
A transformation $\alpha:P\Rightarrow Q$ is a fixed primitive recursive
function $(a,x)\mapsto\alpha_a(x)$ which preserves membership and fibre
equality and satisfies
\begin{equation}
 Q(f)(\alpha_b(x))\approx_a\alpha_a(P(f)x).
 \label{eq:presheaf-naturality}
\end{equation}
Transformations are identified pointwise modulo the target fibre relation.
\end{definition}

There is no assumption that $P(a)$ is a separately existing set object.  The
notation means a predicate on natural numbers with a specified equivalence
relation.  Likewise, writing $P:C^{\mathrm{op}}\to\Set$ is at most a
mnemonic for the data above.  The construction never invokes a universe
containing all sets or all maps.

For an object $A$, the representable presheaf $h_A$ is particularly simple:
\[
 h_A(a)=\Hom_C(a,A),\qquad h_A(f)(u)=u\circ f.
\]
Its membership and equality formulas are the hom formulas of $C$.  Its action
is the fixed composition operation of $C$.  Associativity and identity in
$C$ prove its presheaf laws.  The dependence on the parameter $A$ is uniform:
$A$ is an ordinary numerical parameter, not a code for an arbitrary program.

\section{Yoneda, universal elements, and uniqueness of representations}
\label{sec:yoneda}

\begin{theorem}[Primitive recursive schematic Yoneda lemma]
\label{thm:yoneda-pr}
For every presented presheaf $P$, object $A$, and presented transformation
$\alpha:h_A\Rightarrow P$, define
\[
 \operatorname{ev}_A(\alpha)=\alpha_A(\id_A).
\]
Conversely, for $x\in P(A)$ define a transformation
$\widehat{x}:h_A\Rightarrow P$ by
\begin{equation}
 \widehat{x}_a(f)=P(f)x,\qquad f:a\to A.
 \label{eq:yoneda-construction}
\end{equation}
These constructions respect the relevant equivalence relations and satisfy
\[
 \operatorname{ev}_A(\widehat{x})\approx_A x,
 \qquad
 \widehat{\operatorname{ev}_A(\alpha)}_a(f)
   \approx_a\alpha_a(f).
\]
For each fixed input presentation and transformation term, evaluation is a
primitive recursive term, and reconstruction uses the fixed action term with
the extra numerical parameter $x$.  Evaluation is not a universal evaluator
whose input is the code of an arbitrary program.  Thus the customary notation
$\operatorname{Nat}(h_A,P)\cong P(A)$ has an explicit, first-order schematic
meaning in $\ipra$.
\end{theorem}
\begin{proof}
The identity $\id_A$ is a valid element of $h_A(A)$, so preservation of
membership by $\alpha$ makes its evaluation a valid member of $P(A)$.

For the converse, if $f:a\to A$, action typing gives $P(f)x\in P(a)$.
If $f\sim f'$, action congruence gives $P(f)x\approx_a P(f')x$.
For $k:b\to a$, the required naturality equation is exactly
\[
 P(k)(\widehat{x}_a(f))
 =P(k)(P(f)x)
 \approx_b P(f\circ k)x
 =\widehat{x}_b(h_A(k)(f)).
\]
This is the composition law of the presheaf, used in its symmetric direction.
If $x\approx_A y$, action congruence makes $\widehat{x}$ and $\widehat{y}$
pointwise equal.  If $\alpha$ and $\beta$ are pointwise equal, their values
at $(A,\id_A)$ are equal, so evaluation also respects equality.

The first inverse law is
$\operatorname{ev}_A(\widehat{x})=P(\id_A)x\approx_A x$.
For the second, use naturality of $\alpha$ at $f:a\to A$ and apply it to
$\id_A\in h_A(A)$.  It says
\[
 P(f)(\alpha_A(\id_A))
 \approx_a\alpha_a(h_A(f)(\id_A))
 =\alpha_a(\id_A\circ f)
 \approx_a\alpha_a(f).
\]
The final step is the identity law in $C$, followed by congruence of
$\alpha_a$.  This is exactly the required pointwise equality.  No element
has been chosen by search, and no equality was tested.  The witness term
is the action function of $P$ applied to its two displayed arguments.
\end{proof}

\begin{proposition}[Naturality of Yoneda in both variables]
For $\tau:P\Rightarrow Q$, the Yoneda correspondences commute with applying
$\tau$.  For an arrow $u:A\to B$, they commute with precomposition by the
transformation $h_u:h_A\Rightarrow h_B$, where
$(h_u)_a(f)=u\circ f$, and with the action map $P(u):P(B)\to P(A)$.
\end{proposition}
\begin{proof}
The first assertion at evaluation is literal:
\[
 \operatorname{ev}_A(\tau\alpha)
 =\tau_A(\alpha_A(\id_A)).
\]
At reconstruction it is naturality of $\tau$:
\[
 \tau_a(P(f)x)\approx_a Q(f)(\tau_A(x)).
\]
For the second assertion, let $\alpha:h_B\Rightarrow P$.  Evaluation of
$\alpha h_u$ gives
\[
 \alpha_A(u\circ\id_A)\approx_A\alpha_A(u)
 \approx_A P(u)(\alpha_B(\id_B)),
\]
where the last equality is the second inverse identity already proved.
Equivalently, starting with $x\in P(B)$, reconstruction after applying $P(u)$
has component
\[
 P(f)(P(u)x)\approx_a P(u\circ f)x,
\]
which is reconstruction followed by precomposition with $h_u$.
\end{proof}

\begin{theorem}[The Yoneda embedding is schematically fully faithful]
\label{thm:yoneda-fully-faithful}
The association $A\mapsto h_A$, $u\mapsto h_u$ preserves identities and
composition.  Every fixed presented transformation $\alpha:h_A\Rightarrow h_B$
has the unique inducing arrow, modulo hom equality,
\[
 u=\alpha_A(\id_A):A\to B.
\]
In particular, if $h_u$ and $h_v$ are pointwise equal, then $u\sim v$.
\end{theorem}
\begin{proof}
The component of $h_{\id_A}$ at $f$ is $\id_A\circ f\sim f$.
The component of $h_v h_u$ is $v\circ(u\circ f)\sim(v\circ u)\circ f$,
which is the component of $h_{v\circ u}$.  Thus the association is functorial
in the schematic sense just explained.

Apply Theorem~\ref{thm:yoneda-pr} with $P=h_B$.  The reconstruction of
$u\in h_B(A)$ is
$\widehat{u}_a(f)=h_B(f)(u)=u\circ f$, precisely $h_u$.  Its second inverse
identity gives $\alpha_a(f)\sim u\circ f$.  If also
$\alpha_a(f)\sim v\circ f$ for all $a,f$, put $a=A$ and $f=\id_A$.
The identity laws give $u\sim v$.  This proves fullness and faithfulness
without a general function space or any choice operation.
\end{proof}

\begin{corollary}[Isomorphism is detected by representables]
If presented transformations $\alpha:h_A\Rightarrow h_B$ and
$\beta:h_B\Rightarrow h_A$ are pointwise inverse, the arrows
$u=\alpha_A(\id_A)$ and $v=\beta_B(\id_B)$ satisfy
$v\circ u\sim\id_A$ and $u\circ v\sim\id_B$.
\end{corollary}
\begin{proof}
By fullness, $\alpha=h_u$ and $\beta=h_v$ pointwise.  Functoriality turns
$\beta\alpha=\id$ into $h_{v\circ u}=h_{\id_A}$.  Faithfulness gives
$v\circ u\sim\id_A$.  The other identity is obtained from
$\alpha\beta=\id$ in the same way.  Every arrow used in the conclusion has
already been displayed.
\end{proof}

\subsection{Universal elements and the category of elements}

The category of elements $\int P$ has object codes $(a,x)$ with
$a\in C$ and $x\in P(a)$.  A morphism
\[
 (a,x)\longrightarrow(b,y)
\]
is a morphism $f:a\to b$ satisfying $P(f)y\approx_a x$.
Its equality is the equality of the underlying arrows of $C$.  Pairing and
projection of codes are primitive recursive.  Identity and composition are
inherited from $C$.

\begin{lemma}
The displayed data define an arithmetic presentation of a category.
\end{lemma}
\begin{proof}
For identity, $P(\id_a)x\approx_a x$ is a presheaf law.  For composition,
let $f:(a,x)\to(b,y)$ and $g:(b,y)\to(c,z)$.  Then
\[
 P(g\circ f)z\approx_a P(f)(P(g)z)
 \approx_a P(f)y\approx_a x.
\]
Thus $g\circ f$ has the required element condition.  Congruence of the action
shows that changing an arrow to an equal arrow preserves this condition.
The equivalence, congruence, associativity, and identity laws are inherited
from $C$ with the indicated typing hypotheses.
\end{proof}

\begin{definition}[Chosen universal element]
A chosen universal element of $P$ is a valid pair $(A,x)$ together with a
fixed primitive recursive function $r(a,z)$ such that, whenever $z\in P(a)$,
\[
 r(a,z):a\to A,\qquad P(r(a,z))x\approx_a z,
\]
and every $f:a\to A$ satisfying $P(f)x\approx_a z$ satisfies
$f\sim r(a,z)$.  This is a terminal object of $\int P$ equipped with a
primitive recursive mediator function.
\end{definition}

\begin{theorem}[Representability by a chosen universal element]
A chosen universal element $(A,x,r)$ determines an isomorphism of presented
presheaves $h_A\cong P$.  Conversely, a presented isomorphism
$\theta:h_A\cong P$ determines such a chosen universal element.  The
constructions and their laws are schematic theorems of $\ipra$.
\end{theorem}
\begin{proof}
For the forward direction, put $\theta=\widehat{x}$ and
$\rho_a(z)=r(a,z)$.  The definition gives
$\theta_a(\rho_a(z))\approx_a z$.  It also implies
$\rho_a(\theta_a(f))\sim f$, since $f$ itself satisfies the relevant
universal condition.

We must check that $\rho$ is an extensional natural transformation, rather
than merely a family of inverse values.  If $z\approx_a z'$, then
$P(r(a,z))x\approx_a z'$; uniqueness for $z'$ gives
$r(a,z)\sim r(a,z')$.  Thus $\rho_a$ is extensional.  For $k:b\to a$,
\[
 P(r(a,z)\circ k)x
 \approx_b P(k)(P(r(a,z))x)
 \approx_b P(k)z.
\]
The uniqueness condition at $(b,P(k)z)$ gives
$r(a,z)\circ k\sim r(b,P(k)z)$.  This is naturality of $\rho:P\Rightarrow h_A$.
Therefore $\theta$ and $\rho$ form the required presheaf isomorphism.

Conversely, let $\theta:h_A\Rightarrow P$ and $\rho:P\Rightarrow h_A$ be
presented inverse transformations.  Set $x=\theta_A(\id_A)$ and
$r(a,z)=\rho_a(z)$.  Yoneda gives
$\theta_a(f)\approx_a P(f)x$.  Hence
\[
 P(r(a,z))x\approx_a\theta_a(\rho_a(z))\approx_a z.
\]
If $P(f)x\approx_a z$, the same Yoneda identity gives
$\theta_a(f)\approx_a z$.  Applying the extensional map $\rho_a$ and the
inverse law yields $f\sim\rho_a(z)=r(a,z)$.  This proves uniqueness.
\end{proof}

The adjective ``chosen'' is substantive.  A bare assertion that every object
of $\int P$ has some unique arrow to $(A,x)$ does not, by itself, supply a
single primitive recursive mediator function.  Ordinary intuitionistic logic
allows the corresponding local existential reasoning, but our claim that a
\emph{presented} inverse transformation has been constructed requires the
additional function $r$.  No choice or witness-extraction theorem is being
silently invoked.

\begin{theorem}[Uniqueness of representations, with the canonical arrow]
Suppose $\theta:h_A\cong P$ and $\psi:h_B\cong P$ are presented
isomorphisms.  There is a unique arrow $u:A\to B$, modulo hom equality,
such that
\[
 \psi_a(u\circ f)\approx_a\theta_a(f)
 \qquad(f:a\to A).
\]
It is an isomorphism, and its two directions are explicitly
\[
 u=\psi_A^{-1}(\theta_A(\id_A)),\qquad
 v=\theta_B^{-1}(\psi_B(\id_B)).
\]
\end{theorem}
\begin{proof}
The composite transformation $\psi^{-1}\theta:h_A\Rightarrow h_B$ is
induced by its evaluation at $\id_A$, which is the displayed $u$.
Therefore $\psi^{-1}\theta=h_u$ pointwise.  Composing pointwise with $\psi$
gives the required compatibility.  Any other compatible $u'$ satisfies
$h_{u'}=\psi^{-1}\theta=h_u$, so Yoneda faithfulness gives $u'\sim u$.
The analogous construction gives $h_v=\theta^{-1}\psi$.  These two
transformations are inverse; the preceding corollary gives
$v\circ u\sim\id_A$ and $u\circ v\sim\id_B$.
\end{proof}

\subsection{The covariant form}

For completeness, suppose $Q$ is a presented \emph{covariant} setoid-valued
functor: $f:a\to b$ gives $Q(f):Q(a)\to Q(b)$, with the covariant functor
laws.  Put $k_A(a)=\Hom_C(A,a)$.  The same statement holds with
$k_A$ in place of $h_A$:
\[
 \alpha\longmapsto\alpha_A(\id_A),\qquad
 x\longmapsto\bigl(f\longmapsto Q(f)x\bigr).
\]
To verify the nontrivial inverse directly, naturality at $f:A\to a$ gives
$Q(f)(\alpha_A(\id_A))\approx_a\alpha_a(f\circ\id_A)
\approx_a\alpha_a(f)$.  Naturality of the reconstruction is the covariant
composition law.  A universal element in this variance is a chosen
\emph{initial} object in the covariant category of elements, whose arrow
condition is $Q(f)x\approx_b y$.  Thus both variances are covered by explicit
formulas rather than an appeal to an unconstructed opposite functor category.

\section{Density and co-Yoneda by explicit normalization}
\label{sec:density}

There is a useful strengthening of the preceding analysis.  The density
formula for a presheaf can be proved in a form that avoids assuming the
existence of arbitrary colimits.  The key is that each apparent quotient
class has a canonical representative computable by one presheaf action.

Fix $b\in C$.  A generator over $b$ is a triple
\[
 (a,x,f),\qquad x\in P(a),\quad f:b\to a.
\]
Define its value by $\nu_b(a,x,f)=P(f)x\in P(b)$.  Put
\begin{equation}
 (a,x,f)\equiv_b(a',x',f')
 \quad\Longleftrightarrow\quad
 P(f)x\approx_b P(f')x'.
 \label{eq:density-equality}
\end{equation}
This is an arithmetic equivalence relation; it is not implemented by a
search for an equivalence class.  Let $D_P(b)$ denote this presented setoid.
An arrow $k:c\to b$ acts by
\[
 D_P(k)(a,x,f)=(a,x,f\circ k).
\]
The presheaf action law shows that this action respects
\eqref{eq:density-equality}, and the category laws prove its identity and
composition laws.  Thus $D_P$ is a presented presheaf.

\begin{theorem}[Explicit co-Yoneda normal form]
\label{thm:co-yoneda-normal}
There is a presented natural isomorphism $D_P\cong P$ with maps
\[
 \nu_b(a,x,f)=P(f)x,\qquad
 \iota_b(z)=(b,z,\id_b).
\]
Moreover, the relation \eqref{eq:density-equality} is exactly the least
setoid equivalence compatible with fibre and hom equalities and the elementary
identifications
\begin{equation}
 (a,x,f)\equiv_b(a',x',u\circ f)
 \quad\text{when }u:(a,x)\to(a',x')\text{ in }\int P.
 \label{eq:density-elementary}
\end{equation}
Here ``least'' is a schematic implication about any fixed arithmetic
relation satisfying these closure conditions; it does not quantify over a
set of all relations.
\end{theorem}
\begin{proof}
The maps have the required types.  The composite $\nu_b\iota_b$ sends $z$
to $P(\id_b)z\approx_b z$.  The other composite sends $(a,x,f)$ to
$(b,P(f)x,\id_b)$, whose value is
$P(\id_b)(P(f)x)\approx_b P(f)x$.  The definition of $\equiv_b$ therefore
makes this composite the identity modulo the setoid relation.

Naturality of $\nu$ is the identity
$P(f\circ k)x\approx_c P(k)(P(f)x)$.  For $\iota$, the images
$D_P(k)(b,z,\id_b)=(b,z,\id_b\circ k)$ and
$(c,P(k)z,\id_c)$ have equal values under $\nu_c$ by the identity and
composition laws; hence they are $\equiv_c$-equal.  This proves naturality.

We next identify the quotient relation intrinsically.  Fibre equality and
arrow equality plainly preserve values.  If $u:(a,x)\to(a',x')$, its
defining condition says $P(u)x'\approx_a x$.  Consequently
\[
 P(u\circ f)x'\approx_b P(f)(P(u)x')\approx_b P(f)x.
\]
Thus every elementary identification preserves $\equiv_b$.

Conversely, let $R$ be any equivalence relation on generators that respects
the stated fibre equalities and contains the elementary identifications.
For each generator $t=(a,x,f)$, the arrow
\[
 f:(b,P(f)x)\longrightarrow(a,x)
\]
is an arrow of $\int P$ by reflexivity.  Applying the elementary relation
to this arrow and $\id_b$ gives
\[
 (b,P(f)x,\id_b)\ R\ (a,x,f\circ\id_b)\ R\ (a,x,f),
\]
where the last step uses the specified compatibility with hom equality.
If $t\equiv_b t'$, their two values are fibre-equal, so the corresponding
canonical triples are $R$-equal.  Transitivity and symmetry now give $t R t'$.
This uses only a fixed finite sequence of equalities, regardless of the
numerical codes.  In particular, there is no appeal to induction on an
arbitrarily long chain of quotient generators.
\end{proof}

The customary coend notation for this result is
\[
 P(b)\cong\int^{a\in C}P(a)\times\Hom_C(b,a).
\]
In the present foundation, this expression means precisely the coded
setoid and normalization maps above.  It should not be read as a theorem
asserting all coends exist in a global category of sets.

\begin{theorem}[Schematic density as a universal cocone]
\label{thm:density-cocone}
For $(a,x)\in\int P$, let $j_{a,x}:h_a\Rightarrow P$ be the Yoneda
transformation $j_{a,x,b}(f)=P(f)x$.  This family is a cocone over the diagram
$(a,x)\mapsto h_a$, with an arrow $u$ acting by postcomposition $h_u$.
It has the following universal property.  For every presented presheaf $Q$
and every \emph{uniformly primitive recursive} family of transformations
$\lambda_{a,x}:h_a\Rightarrow Q$ satisfying
\begin{equation}
 \lambda_{a,x,b}(f)\approx_b
 \lambda_{a',x',b}(u\circ f)
 \quad\bigl(u:(a,x)\to(a',x')\bigr),
 \label{eq:density-cocone}
\end{equation}
there is a unique presented transformation $\delta:P\Rightarrow Q$ such
that $\delta j_{a,x}=\lambda_{a,x}$ pointwise.  It is explicitly
\begin{equation}
 \delta_b(z)=\lambda_{b,z,b}(\id_b).
 \label{eq:density-mediator}
\end{equation}
\end{theorem}
\begin{proof}
The cocone law for $j$ says
$P(f)x\approx_b P(u\circ f)x'$, and this was proved in the preceding theorem.
The displayed formula for $\delta$ is primitive recursive because the family
$\lambda$ is assumed uniform in all of its numerical arguments.  Its value
belongs to $Q(b)$ by typing of $\lambda_{b,z}$.

To check extensionality, suppose $z\approx_b z'$.  The identity arrow
$\id_b$ is an arrow $(b,z)\to(b,z')$ in $\int P$, since
$P(\id_b)z'\approx_b z$.  Applying the cocone equation at $\id_b$ gives
\[
 \lambda_{b,z,b}(\id_b)\approx_b
 \lambda_{b,z',b}(\id_b\circ\id_b)
 \approx_b\lambda_{b,z',b}(\id_b).
\]
The last equality uses the identity law and extensionality of
$\lambda_{b,z',b}$.  This is $\delta_b(z)\approx_b\delta_b(z')$.

For naturality, let $k:c\to b$.  The arrow
$k:(c,P(k)z)\to(b,z)$ lies in $\int P$.  The cocone condition gives
\[
 \delta_c(P(k)z)=\lambda_{c,P(k)z,c}(\id_c)
 \approx_c\lambda_{b,z,c}(k).
\]
Naturality of the fixed transformation $\lambda_{b,z}:h_b\Rightarrow Q$
at $k$, evaluated at $\id_b$, gives
\[
 \lambda_{b,z,c}(k)\approx_c
 Q(k)(\lambda_{b,z,b}(\id_b))=Q(k)(\delta_b(z)).
\]
These two equalities prove naturality of $\delta$.

To prove factorization, take $f:b\to a$ and $x\in P(a)$.  The arrow
$f:(b,P(f)x)\to(a,x)$ belongs to $\int P$, so the cocone condition gives
\[
 \delta_b(j_{a,x,b}(f))
 =\lambda_{b,P(f)x,b}(\id_b)
 \approx_b\lambda_{a,x,b}(f).
\]
Finally, let $\delta'$ be another factorization.  Applying its factorization
equation to the canonical generator $(b,z,\id_b)$ yields
\[
 \delta'_b(P(\id_b)z)\approx_b\lambda_{b,z,b}(\id_b)=\delta_b(z).
\]
The presheaf identity law and extensionality of $\delta'_b$ replace the
left side by $\delta'_b(z)$.  Thus $\delta'$ equals $\delta$ pointwise.
\end{proof}

\paragraph{Foundational content.}
The density theorem has an unbounded object predicate, but its proof still
requires no unbounded search, induction on that predicate, choice of
representatives, or generic recursion over formulas.  It constructs the
mediator by a single substitution into the \emph{supplied uniform} cocone
program.  Removing that uniformity assumption would remove the reason that
\eqref{eq:density-mediator} is primitive recursive.  This distinction is
precisely what allows a substantial part of presheaf theory to survive in
$\ipra$ while avoiding a claim that all classical presheaf categories have
been internalized.

\section{Universal constructions with explicit operations}\label{sec:limits}

In this section, an equality of parallel arrows always means the chosen
hom-setoid equality.  We suppress its source and target parameters, but never
identify two object codes merely because their objects are isomorphic.  Every
operation mentioned below is a fixed primitive recursive operation on codes.
Its typing and algebraic laws are hypotheses already proved for the
presentation under consideration.  An operation may return arbitrary junk on
inputs outside its stated domain.

\subsection{Terminal objects, products, and uniqueness}

A \emph{chosen terminal object} consists of an object $1$ and an operation
$!_a:a\to1$ such that every arrow $u:a\to1$ equals $!_a$.  A chosen initial
object is the dual construction.  For an initial object $0$, its unique arrow
$0\to a$ is supplied as an operation; the statement that an arrow exists for
each $a$ is not silently converted into such an operation.

A \emph{chosen binary product} of $a,b$ consists of an object $a\times b$,
projections $p:a\times b\to a$ and $q:a\times b\to b$, and a pairing operation
\[
 (u:x\to a,\ v:x\to b)\longmapsto\langle u,v\rangle:x\to a\times b
\]
satisfying
\[
 p\langle u,v\rangle=u,\qquad q\langle u,v\rangle=v,\qquad
 \langle pw,qw\rangle=w.
\]
Composition is written by juxtaposition, with the arrow on the right applied
first.  These three equations include both existence and uniqueness.  They
also prove that pairing respects equality: if $u=u'$ and $v=v'$, the two
pairings have equal projections and hence are equal.

\begin{proposition}[Uniqueness of universal objects]\label{prop:universal-unique}
Two terminal objects are uniquely isomorphic.  Two products of the same
ordered pair are uniquely isomorphic by an isomorphism preserving projections.
More generally, two limits of a fixed diagram are uniquely isomorphic by an
isomorphism preserving their cones.  All these isomorphisms and their inverses
are obtained from the supplied universal operations by composition.
\end{proposition}
\begin{proof}
For terminal objects $t,t'$, let $u:t\to t'$ and $v:t'\to t$ be their supplied
unique arrows.  Since $vu$ and $\id_t$ both map $t$ to $t$, uniqueness at $t$
gives $vu=\id_t$; similarly $uv=\id_{t'}$.  Any arrow $t\to t'$ equals $u$.

For products $(p:P\to a,q:P\to b)$ and $(p':P'\to a,q':P'\to b)$, put
$u=\langle p,q\rangle':P\to P'$ and
$v=\langle p',q'\rangle:P'\to P$.  Then $pvu=p'u=p$ and $qvu=q'u=q$, so
product uniqueness gives $vu=\id_P$.  The other composite is analogous.
An arrow preserving both projections must equal the indicated pairing.

For limits, let $u:L\to L'$ be the mediating arrow from the cone of $L$
into the limiting cone of $L'$, and define $v:L'\to L$ conversely.  Every
projection of $vu$ equals the corresponding projection of $\id_L$; limit
uniqueness gives their equality.  Interchanging the two cones proves the
other composite equation and the same uniqueness argument proves the final
claim.  These are finite derivations in intuitionistic first-order logic:
neither excluded middle nor induction is used.
\end{proof}

An \emph{equalizer} of $f,g:a\to b$ is an object $E$ and $e:E\to a$ with
$fe=ge$, together with a supplied operation taking $u:x\to a$ and the premise
$fu=gu$ to an arrow $\operatorname{lift}(u):x\to E$.  Its laws are
\[
 e\operatorname{lift}(u)=u,
 \qquad ev=u\ \Longrightarrow\ v=\operatorname{lift}(u).
\]
The operation need not test the premise: its correctness is asserted only
under that premise.  Equalizers are monomorphisms: if $er=es$, both $r$ and
$s$ are mediating arrows for the same arrow $er$, so they are equal.

A \emph{pullback} of $f:a\to c$ and $g:b\to c$ consists of
$p:P\to a$, $q:P\to b$ with $fp=gq$, and a supplied operation
$(u,v)\mapsto\langle u,v\rangle_{f,g}$ for compatible pairs $fu=gv$, having
those projections and uniquely determined by them.

\begin{proposition}[Pullbacks from products and equalizers]
If the relevant chosen product and equalizer are supplied, the pullback of
$f:a\to c$ and $g:b\to c$ is obtained by equalizing
\[
 f p_a,\ g p_b:a\times b\longrightarrow c.
\]
Conversely, a terminal object and chosen pullbacks produce binary products;
chosen binary products and pullbacks produce equalizers.
\end{proposition}
\begin{proof}
Let $e:P\to a\times b$ equalize $fp_a,gp_b$ and give $P$ the projections
$p_a e,p_b e$.  Their compatibility is the equalizer equation.  If $fu=gv$,
then $\langle u,v\rangle$ equalizes the two arrows, so its equalizer lift has
the required projections.  An arrow with those projections satisfies
$ew=\langle u,v\rangle$ by product uniqueness; equalizer uniqueness then
identifies it with the lift.

The pullback of the two unique arrows $a\to1\leftarrow b$ is a product,
because every pair $u:x\to a,v:x\to b$ is compatible by terminal uniqueness.
For the last assertion, pull back $(f,g):a\to b\times b$ along the diagonal
$\Delta_b=\langle\id_b,\id_b\rangle:b\to b\times b$.  Write the resulting
projections as $e:E\to a$ and $r:E\to b$.  The pullback equation gives
$fe=r=ge$.  If $fu=gu$, the pair $(u,fu)$ is compatible and therefore has a
unique lift $w$ with $ew=u$.  Any other $w'$ with $ew'=u$ has
$rw'=few'=fu$, so pullback uniqueness proves $w'=w$.
\end{proof}

\subsection{All fixed finite limits}

Fix externally a finite category $J$: a finite list of objects, a finite list
of arrows with specified endpoints, and fixed identity and composition
tables.  For a diagram $D:J\to C$, a cone with vertex $x$ is a finite tuple
$c_i:x\to D_i$ satisfying $D(u)c_i=c_j$ for every listed arrow $u:i\to j$.
A chosen limit consists of a cone $p_i:L\to D_i$ and a primitive recursive
operation $m(x,(c_i))$ such that $p_i m=c_i$, and such that any arrow with
those equations equals $m$.  A cone predicate is an arithmetical formula; it
need not be decidable.  The operation is only required to have its stated
properties under that predicate.

\begin{theorem}[Finite-limit construction]\label{thm:finite-limits}
For each fixed finite category $J$, chosen terminal objects, binary products,
and equalizers yield a chosen limit for every $J$-diagram.  The constructor
and the mediating operation are primitive recursive compositions of the
supplied operations.  The proof translates into \ipra{} without any induction
on a formula containing the category's equality predicate.
\end{theorem}
\begin{proof}
Form the finite products
\[
 A=\prod_{i\in\operatorname{Ob}J}D_i,
 \qquad B=\prod_{(u:i\to j)\in\operatorname{Arr}J}D_j.
\]
They mean fixed parenthesized iterations of binary products; an empty product
means the terminal object.  Write $\pi_i:A\to D_i$ and
$\rho_u:B\to D_j$.  By repeated pairing define $s,t:A\to B$ by
\[
 \rho_u s=D(u)\pi_i,\qquad \rho_u t=\pi_j.
\]
Take their equalizer $e:L\to A$ and put $p_i=\pi_i e$.  Then
$D(u)p_i=D(u)\pi_i e=\rho_u se=\rho_u te=\pi_j e=p_j$, so these arrows
form a cone.

Given a cone $(c_i)$, repeated pairing supplies $c:x\to A$ with
$\pi_i c=c_i$.  For every arrow $u:i\to j$,
$\rho_u sc=D(u)c_i=c_j=\rho_u tc$.  Product uniqueness gives $sc=tc$.
Consequently its equalizer lift $m:x\to L$ exists by the supplied operation
and satisfies $p_i m=\pi_i em=\pi_i c=c_i$.  If $v:x\to L$ has the same
cone projections, then all product projections of $ev$ and $c$ agree, so
$ev=c$.  Equalizer uniqueness yields $v=m$.

All lists in this proof have an externally fixed length.  Thus each use of
``for every listed arrow'' expands to a finite conjunction, and repeated
pairing expands to a fixed term.  The argument is a finite intuitionistic
proof for that shape.  It does not infer a single internal theorem for
arbitrarily long diagrams by induction on an undecidable predicate.  A
uniform varying-length version requires an additional verified coding and
invariant argument, as discussed in the foundational section.
\end{proof}

Applying the same argument in $C^{\op}$ proves the dual theorem: a chosen
initial object, binary coproducts, and coequalizers give all fixed finite
colimits.  In particular, a pushout is the coequalizer of the two induced
arrows into a coproduct.  The hypotheses supply the coequalizer operation;
this theorem does not manufacture quotient objects in an arbitrary category.

\subsection{Idempotents and their canonical completion}

An idempotent is an endomorphism $e:a\to a$ satisfying $e^2=e$.  A
\emph{splitting} consists of $r:a\to b$ and $s:b\to a$ satisfying
$sr=e$ and $rs=\id_b$.  A splitting is therefore a displayed retraction,
not an unspecified existential witness.

\begin{theorem}[Explicit idempotent completion]\label{thm:karoubi}
Every presented category $C$ has a presented category $\operatorname{Kar}(C)$
whose objects are pairs $(a,e)$ with $e^2=e$, and whose arrows
$(a,e)\to(b,d)$ are arrows $f:a\to b$ satisfying
\[
 fe=f=df.
\]
Composition is inherited from $C$, and the identity of $(a,e)$ is $e$.
The inclusion $a\mapsto(a,\id_a)$ is fully faithful with the identity
operation on hom-codes.  Every idempotent in $\operatorname{Kar}(C)$ has a
canonical primitive recursive splitting.
\end{theorem}
\begin{proof}
If $f:(a,e)\to(b,d)$ and $g:(b,d)\to(c,k)$, then
$(gf)e=g(fe)=gf$ and $k(gf)=(kg)f=gf$, so $gf$ is a valid arrow.
The equations $fe=f=df$ make $e,d$ the appropriate identities;
associativity follows from that in $C$.  For identity idempotents the support
conditions impose no extra restriction, proving full faithfulness.

Let $p:(a,e)\to(a,e)$ satisfy $p^2=p$.  The object $(a,p)$ is valid.
Use the underlying arrow $p$ for both
$r:(a,e)\to(a,p)$ and $s:(a,p)\to(a,e)$.  Their support equations follow
from $pe=p=ep$ and $p^2=p$.  We have $sr=p$ as an endomorphism of $(a,e)$,
while $rs=p=\id_{(a,p)}$.  Pairing and the projection operations on codes
provide the asserted constructors.
\end{proof}

For completeness, suppose $F:C\to D$ is supplied and that every $F(e)$ has
a chosen splitting $F(a)\xrightarrow{r_e}B_e\xrightarrow{s_e}F(a)$,
primitive recursively in the displayed idempotent.  Then
\[
 \widehat F(a,e)=B_e,\qquad
 \widehat F(f)=r_d F(f)s_e
\]
defines a functor.  Indeed,
$r_eF(e)s_e=r_es_er_es_e=\id_{B_e}$, and
\[
 (r_kF(g)s_d)(r_dF(f)s_e)
 =r_kF(g)F(d)F(f)s_e
 =r_kF(gf)s_e.
\]
For the inclusion $I:C\to\operatorname{Kar}(C)$, this is an extension
\emph{up to a specified natural isomorphism}, not necessarily a literal
identity of object codes.  Its component
$s_{\id_a}:\widehat F(Ia)\to Fa$ has inverse $r_{\id_a}$, and
\[
 s_{\id_b}\widehat F(I f)
 =s_{\id_b}r_{\id_b}F(f)s_{\id_a}
 =F(f)s_{\id_a}.
\]
Here $s_{\id_b}r_{\id_b}=F(\id_b)=\id_{Fb}$ modulo the given hom equality.
If the identity splittings are chosen as the canonical trivial splittings,
the object maps agree literally.  The theorem does not assert that arbitrary
idempotents already split in $C$ or $D$.

\section{A uniform finite universe, including all empty cases}\label{sec:finite}

The preceding structural theory applies to infinite presentations. A second
useful result is more algorithmic: the entire finite theory can be made
uniform in numerical table codes. Here every verification quantifier is
bounded, and every loop invariant has a PR Boolean characteristic function.
This section supplies the constructions and their correctness proofs, rather
than assuming that an unspecified universal object can be found.

\subsection{The category of all coded finite sets}
Write $[n]=\{0,\ldots,n-1\}$, allowing $n=0$. Define
\[
 b=m+1,\qquad d_b(c,i)=\bigl\lfloor c/b^i\bigr\rfloor\bmod b.
\]
The divisor is positive, including when $m=0$. A table code for $[n]\to[m]$
is a number $c$ satisfying
\[
 c<(m+1)^n,\qquad\forall i<n\; d_{m+1}(c,i)<m.                 \tag{FT}
\]
The bounded formula (FT) is a PR predicate. Its interpretation is
$i\mapsto d_{m+1}(c,i)$. Conversely a bounded list $v_i<m$ is encoded by
$\sum_{i<n}v_i(m+1)^i$. The digit and reconstruction equations follow by
quantifier-free induction on the coded length, using division with remainder.
For $m=0$, a valid table exists exactly when $n=0$, and then its only code
is $0$. This convention eliminates a common hidden positivity assumption.

Identity is the code of $i\mapsto i$. For $f:[n]\to[m]$ and $g:[m]\to[k]$,
composition is the code of
\[
 i\longmapsto d_{k+1}\bigl(g,d_{m+1}(f,i)\bigr).
\]
The displayed bound on the inner digit gives the typing needed for the
outer table. Associativity follows at each input by associativity of
substitution of three table lookups, and table reconstruction then gives
literal equality of the codes. The same pointwise proof gives both unit
laws. The object predicate is true of every natural number; equality of
arrows is numerical equality. Hence this is a uniform iPRA presentation
$\mathsf{Fin}$, not a separate category chosen for each size bound.

The elementary rank/select tool will recur. Given a bit test $K(i)$ and
bound $B$, let
\[
 r(t)=\sum_{i<t}K(i),\quad \ell=r(B),\quad
 s(j)=\min\{i<B:K(i)=1\text{ and }r(i)=j\},                  \tag{RS}
\]
where the minimum returns $B$ if no witness exists. Then
\[
 \begin{split}
 j<\ell&\Rightarrow s(j)<B,\ K(s(j))=1,\ r(s(j))=j,\\
 i<B\wedge K(i)=1&\Rightarrow r(i)<\ell,\ s(r(i))=i.
 \end{split}                                               \tag{RS'}
\]
To prove these facts, scan $i=0,\ldots,B-1$: the counter increases by one
exactly at a successful entry. Its successive successful values are
$0,\ldots,\ell-1$, each attained once. The assertion for any current
prefix is a conjunction of bounded bit tests, so this is a legitimate
quantifier-free induction. There is no unbounded minimization.

\subsection{Initial objects, products, equalizers and all finite limits}
The object $[0]$ is initial: a table with empty domain has its unique code
$0$. The object $[1]$ is terminal: every entry must be $0$, so again there
is one table. The product of $[n]$ and $[m]$ is $[nm]$, with coordinates
\[
 \pi_1(t)=\lfloor t/m\rfloor,\qquad \pi_2(t)=t\bmod m.
\]
Division by zero is totalized by zero. If $nm=0$, the domain is empty, so
these defaults never enter a typing assertion. If $t<nm$, then $m>0$,
and the quotient and remainder equations give the required ranges.
Given $f:[z]\to[n]$, $g:[z]\to[m]$, the mediator is
$u(x)=m f(x)+g(x)$. Its projections are $f,g$, and any other mediator
must have the same quotient and remainder, hence the same entries.

For $f,g:[n]\to[m]$, filter $i<n$ by $f(i)=g(i)$. If $\ell$ entries
remain, the equalizer is $s:[\ell]\to[n]$ from (RS). A table
$h:[z]\to[n]$ with $fh=gh$ factors by $u(x)=r(h(x))$.
Equations (RS') prove $su=h$; applying $r$ to this equation proves
uniqueness. When no entry remains, $\ell=0$. A valid cone then has empty
domain, as follows from the same factorization proof.

\begin{theorem}[Uniform finite limits in $\mathsf{Fin}$]\label{thm:finite-limit}
There is a PR procedure which, from a finite table category $J$ and a
table functor $D:J\to\mathsf{Fin}$, constructs a limit, its projections,
and the unique mediating table of any supplied cone. The construction
and its universal property are provable in iPRA. Empty diagrams and
zero-valued objects are included.
\end{theorem}
\begin{proof}
Let the objects of $J$ be $0,\ldots,r-1$, and let $D(i)=[n_i]$. Enumerate
all tuples $x=(x_0,\ldots,x_{r-1})$ with $x_i<n_i$. A convenient
implementation enumerates words in radices $n_i+1$ and filters out the
spare digit in each coordinate. For $r=0$ there is one empty tuple; if
some $n_i=0$ and $r>0$, there are no tuples.
Retain exactly the tuples for which, for every arrow $a:i\to j$ of $J$,
\[
                         D(a)(x_i)=x_j.                     \tag{LC}
\]
This is a finite product of equality bits over the supplied arrow table,
not a parameter whose correctness is assumed. Let $L$ be the increasing
list of retained tuple codes and $\ell$ its length. Projection $p_i$ takes
an index in $L$ to coordinate $i$. Condition (LC) proves the cone laws.

Given a cone $q_i:[z]\to[n_i]$, at each $t<z$ form the tuple
$(q_i(t))_{i<r}$. It satisfies (LC), so (RS') assigns it a unique rank
$u(t)<\ell$. These ranks form the mediating table. Its projections are
$q_i$ by construction. If $v$ also has these projections, the tuple at
$v(t)$ equals the tuple at $u(t)$ coordinatewise; reconstruction and
injectivity of selection yield $v(t)=u(t)$.

Every loop ranges over a displayed finite code bound: objects, arrows,
tuple codes or cone inputs. The invariants say that a retained tuple
satisfies each earlier equation, or that the constructed output entries
have the required ranks and coordinates. All are bounded arithmetic
formulas and are Booleanized as in Lemma~\ref{lem:bounded}. This yields
uniform numerical iPRA proofs, including the input sizes as variables.
For the empty diagram, the sole empty tuple is retained and the limit
is $[1]$. If a diagram object is empty, its full product is empty and
the limit is $[0]$; a cone into it must itself have empty domain.
\end{proof}

\subsection{Finite quotients and all finite colimits}
For two tables $f,g:[p]\to[q]$, start with a label table
$\lambda_0(v)=v$ for $v<q$. At stage $i<p$ let
$a=\lambda_i(f(i))$ and $b=\lambda_i(g(i))$. Replace every label equal
to $b$ by $a$, leaving all other labels unchanged. If $a=b$ this does
nothing. Every label remains below $q$, and the label table remains
idempotent: its labels are precisely fixed points of the table.

The following stronger invariants explain why this merging does not
identify anything unjustified. After stage $i$:
\begin{enumerate}
\item the endpoints of every processed pair have equal labels;
\item if $h:[q]\to[z]$ equalizes all processed pairs, then
      $h(v)=h(\lambda_i(v))$ for every $v<q$;
\item each $v$ is connected to $\lambda_i(v)$ by a chain of processed
      pair identifications, used in either direction.
\end{enumerate}
For (2), suppose the next pair is equalized. The old invariant gives
$h(a)=h(f(i))=h(g(i))=h(b)$; replacing label $b$ by $a$ therefore
preserves $h(v)=h(\lambda(v))$. For (3), concatenate the two old paths
from $f(i),g(i)$ to $a,b$ with the new pair, and attach this path when
a label changes. If the old paths have length at most $B_i$, the path from $b$ to $a$
has length at most $2B_i+1$, and an updated path has length at most
$3B_i+1$. Thus $B_0=0$, $B_{i+1}=3B_i+1$ is an explicit PR bound.
A coded path table can be updated by the same primitive recursion.
For factorization only (2) is needed. Invariant (3), together with
(1) applied along a coded path, additionally proves that equality of
final labels is exactly the finite-path relation generated by the
pairs. This is a claim about a coded PR graph and its paths, not a
minimality assertion quantified over arbitrary arithmetic relations.

For invariant (2), fix the target size $z$ and candidate table $h$ as
numerical parameters during induction. The processed-pair condition and
the invariant both quantify only over bounded indices in coded tables.
Their Booleanized implication is therefore quantifier-free; only after
the induction do we generalize over $h,z$.

After all pairs, enumerate the fixed-point labels in increasing order,
say $e_0,\ldots,e_{k-1}$. Define
\[
 Q(v)=\text{the rank of }\lambda_p(v)\text{ among these labels}.
\]
For an equalizing table $h$ set $\bar h(j)=h(e_j)$. Invariant (2)
proves $\bar h Q=h$. Also $Q(e_j)=j$, so any factor has value $h(e_j)$
at $j$; this proves uniqueness. Thus $Q:[q]\to[k]$ is a coequalizer.
If $q=0$, validity of $f,g$ forces $p=0$; the empty label list gives
$k=0$, with exactly the same universal property. All invariants used
to obtain the finite-set universal property concern bounded tables
and their entries, so the induction remains quantifier-free after
Booleanization. It is not induction over arbitrary target categories.

\begin{theorem}[Uniform finite colimits]\label{thm:finite-colimit}
Every finite table diagram in $\mathsf{Fin}$ has a colimit constructed
by a PR algorithm, with a PR operation on cocones returning the unique
mediator. Its iPRA correctness includes empty diagrams and empty sets.
\end{theorem}
\begin{proof}
For $D:J\to\mathsf{Fin}$ form the tagged disjoint union
$U=\coprod_{i\in J}[n_i]$, a finite ordinal obtained by cumulative sums.
For each arrow $a:i\to j$ and $x<n_i$, create the pair
$(i,x),(j,D(a)x)$. The total number of pairs is
$\sum_{a\in\operatorname{Arr}J}n_{\operatorname{src}a}$.
Apply the preceding merging algorithm. Its quotient maps composed with
the tag injections form a cocone. A cocone $q_i:[n_i]\to[z]$
defines $h(i,x)=q_i(x)$ on $U$ and equalizes each generating pair.
The coequalizer factorization therefore gives its unique mediator.
For the empty diagram, $U$ is empty and the colimit is $[0]$.
\end{proof}

\subsection{Exponentials with \texorpdfstring{$0^0=1$}{0 to power 0 equals 1}}
Let $V_{n,m}$ be the increasing list of all codes satisfying (FT), and
write $q=|V_{n,m}|$. There is one choice for an empty list of entries;
appending one entry gives exactly $m$ choices. Therefore the counting
recurrence is $q(0,m)=1$, $q(n+1,m)=m q(n,m)$, whence
$q(n,m)=m^n$ by quantifier-free induction. In particular $0^0=1$,
while $0^n=0$ for $n>0$.

\begin{theorem}[Cartesian closure of coded finite sets]\label{thm:finite-ccc}
In $\mathsf{Fin}$ the object $[q]$, $q=m^n$, is an exponential
$[m]^{[n]}$. Evaluation and currying are uniformly PR, and their
universal equations have iPRA proofs for all $n,m,z$, including zero.
\end{theorem}
\begin{proof}
Let $s$ select a raw code from $V_{n,m}$ and $r$ rank a valid code.
Define $\operatorname{ev}(j,x)=d_{m+1}(s(j),x)$ for $j<q,x<n$.
This is a table on $[q]\times[n]$; a pair never exists if one factor
is empty. For $h:[z]\times[n]\to[m]$, form the row code
\[
 c_t=\sum_{x<n}h(t,x)(m+1)^x,
 \qquad\operatorname{curry}(h)(t)=r(c_t).
\]
Every row is valid, so the rank is below $q$. Digit reconstruction gives
$\operatorname{ev}(\operatorname{curry}(h)(t),x)=h(t,x)$.
Conversely, if $u:[z]\to[q]$, the row obtained by evaluating $u(t)$
is the raw table $s(u(t))$, whose rank is $u(t)$. Thus
\[
 \operatorname{ev}(\operatorname{curry}h\times\id)=h,
 \qquad\operatorname{curry}(\operatorname{ev}(u\times\id))=u.
\]
The first is the factorization law and the second proves uniqueness.
If $n=0$, every row is empty and $q=1$. If $m=0<n$, then $q=0$;
a table $h:[z]\times[n]\to[0]$ forces $z=0$, so its transpose
has the required empty domain. No positivity premise has been omitted.
\end{proof}

\subsection{Finite functor categories and genuine finite Yoneda bijections}
A finite table category consists of finite object and arrow tables,
source, target, identity and composition entries, with a bounded checker
for all typing, unit and associativity equations. Given such $C,D$,
enumerate all assignments on their finite object and arrow domains and
retain precisely those satisfying the functor laws. Likewise enumerate
all object-indexed arrow tuples and retain the natural transformations.
Pointwise identities and vertical composition make a finite table category
$[C,D]$. The verification is the same finite naturality calculation as
in Section~\ref{sec:yoneda}, now with bounded quantifiers on actual tables.

For a finite $\mathsf{Fin}$-valued presheaf $P$, whose fibres are
finite ordinals with numerical equality, all transformations
$h_A\Rightarrow P$ can be enumerated in this manner. The Yoneda
constructions become reciprocal numerical maps between two finite lists.
Consequently $|\operatorname{Nat}(h_A,P)|=|P(A)|$, with an explicit
bijection, rather than merely a schematic correspondence.

Finite adjunction existence is also decidable by bounded enumeration:
enumerate candidate right-adjoint functor tables, then unit and counit
tables, and check naturality and the two triangles. This does not assert
that a right adjoint always exists. It returns either displayed adjunction
data or a finite certificate that no candidate table satisfies the laws.
The bounds may be enormous; primitive recursive arithmetic imposes no
polynomial-time requirement.

\section{Adjunctions, equivalences, and preservation}\label{sec:adjunctions}

\subsection{Two equivalent descriptions of an adjunction}

Let $F:C\to D$ and $G:D\to C$ be fixed presented functors.  A
\emph{displayed hom-adjunction} consists of primitive recursive operations
\[
 \Phi_{a,b}:\Hom_D(Fa,b)\longrightarrow\Hom_C(a,Gb),\qquad
 \Psi_{a,b}:\Hom_C(a,Gb)\longrightarrow\Hom_D(Fa,b),
\]
respecting equality, inverse modulo equality, and natural in the two
variables.  For $u:a'\to a$, $v:b\to b'$, and $f:Fa\to b$, naturality means
\[
 \Phi_{a',b'}(v f F(u))=G(v)\Phi_{a,b}(f)u.
\]
The displayed hom-sets are notation for typed codes modulo the chosen
arithmetic relation.  No set of all maps between them is postulated.

\begin{theorem}[Hom-adjunction and triangle identities]\label{thm:adjunction}
A displayed hom-adjunction is equivalently given by natural transformations
\[
 \eta:1_C\Longrightarrow GF,\qquad
 \varepsilon:FG\Longrightarrow1_D
\]
with primitive recursive component operations, satisfying
\[
 \varepsilon_{Fa}F(\eta_a)=\id_{Fa},\qquad
 G(\varepsilon_b)\eta_{Gb}=\id_{Gb}.
\]
The conversion in each direction is primitive recursive and its correctness
uses only finite intuitionistic equational reasoning.
\end{theorem}
\begin{proof}
Starting with the unit and counit, define
\[
 \Phi(f)=G(f)\eta_a,\qquad \Psi(g)=\varepsilon_bF(g).
\]
Both terms have the stated types and respect equality.  Naturality of the
counit at $f:Fa\to b$, followed by the first triangle, gives
\[
 \Psi\Phi(f)
 =\varepsilon_bFG(f)F(\eta_a)
 =f\varepsilon_{Fa}F(\eta_a)=f.
\]
Naturality of the unit at $g:a\to Gb$, followed by the second triangle,
gives
\[
 \Phi\Psi(g)
 =G(\varepsilon_b)GF(g)\eta_a
 =G(\varepsilon_b)\eta_{Gb}g=g.
\]
For naturality in the source, use $GF(u)\eta_{a'}=\eta_a u$:
\[
 \Phi(fF(u))=G(f)GF(u)\eta_{a'}=G(f)\eta_a u=\Phi(f)u.
\]
Naturality in the target is
$\Phi(vf)=G(v)G(f)\eta_a=G(v)\Phi(f)$.  Combining these proves the stated
bivariable law.

Conversely, let $\Phi,\Psi$ be inverse natural correspondences and set
\[
 \eta_a=\Phi_{a,Fa}(\id_{Fa}),\qquad
 \varepsilon_b=\Psi_{Gb,b}(\id_{Gb}).
\]
Target naturality immediately gives
$\Phi(f)=G(f)\eta_a$.  The inverse correspondence is also natural: for
example, applying $\Phi$ to $\Psi(g)F(u)$ gives $gu$, so injectivity of
$\Phi$ gives $\Psi(gu)=\Psi(g)F(u)$.  The same argument in the target gives
$\Psi(G(v)g)=v\Psi(g)$.  Taking $g=\id_{Gb}$ in the source formula proves
$\Psi(g)=\varepsilon_bF(g)$ for arbitrary $g:a\to Gb$.

The inverse equations at the relevant identities now read
\[
 \varepsilon_{Fa}F(\eta_a)
 =\Psi(\eta_a)=\id_{Fa},\qquad
 G(\varepsilon_b)\eta_{Gb}
 =\Phi(\varepsilon_b)=\id_{Gb}.
\]
For $u:a\to a'$, source naturality and the formula for $\Phi$ yield
$\eta_{a'}u=\Phi(F(u))=GF(u)\eta_a$.  For $v:b\to b'$, target naturality
of $\Psi$ and its formula give
$v\varepsilon_b=\Psi(G(v))=\varepsilon_{b'}FG(v)$.  These are precisely
the naturality equations for the unit and counit.  The two conversion
procedures recover the original data modulo their appropriate equalities.
\end{proof}

This proof explains why existence claims need care.  A family of statements
$\forall a,b\ \exists$ a bijection between hom-collections does not itself
supply the primitive recursive operations $\Phi,\Psi$.  Conversely, once
those operations are supplied, no choice principle is needed in the theorem.

\begin{proposition}[Uniqueness of a specified right adjoint]
For fixed $F$, any two displayed right adjoints $G,G'$ are naturally
isomorphic by explicit component operations.
\end{proposition}
\begin{proof}
Use primes for the second correspondence and counit.  Define
\[
 u_b=\Phi'_{Gb,b}(\varepsilon_b):Gb\to G'b,
 \qquad
 v_b=\Phi_{G'b,b}(\varepsilon'_b):G'b\to Gb.
\]
Then $\Psi'(u_b)=\varepsilon_b$ and $\Psi(v_b)=\varepsilon'_b$.
Naturality of $\Psi$ in its source gives
\[
 \Psi(v_bu_b)=\Psi(v_b)F(u_b)
 =\varepsilon'_bF(u_b)=\Psi'(u_b)=\varepsilon_b
 =\Psi(\id_{Gb}).
\]
Injectivity of $\Psi$ proves $v_bu_b=\id_{Gb}$; the other composite follows
by interchanging primes.  For $h:b\to b'$, applying $\Psi'$ to
$G'(h)u_b$ gives $h\varepsilon_b$.  Applying it to $u_{b'}G(h)$ gives
$\varepsilon_{b'}FG(h)$, which is the same arrow by naturality of
$\varepsilon$.  Thus $G'(h)u_b=u_{b'}G(h)$, proving naturality.
\end{proof}

\subsection{Adjoints preserve supplied universal constructions}

\begin{theorem}[Preservation of limits and colimits]\label{thm:adjoint-preserve}
A displayed right adjoint preserves every supplied fixed finite limit, with
an explicitly constructed primitive recursive mediating operation.  A
displayed left adjoint preserves every supplied fixed finite colimit.
\end{theorem}
\begin{proof}
Suppose $(p_j:l\to H_j)$ is a chosen limit of a fixed finite diagram
$H:J\to D$.  The arrows $G(p_j):Gl\to GH_j$ form a cone by functoriality.
Given a cone $c_j:x\to GH_j$, transpose its components to
$\Psi(c_j):Fx\to H_j$.  For $h:i\to j$ in $J$, naturality gives
\[
 H(h)\Psi(c_i)=\Psi(GH(h)c_i)=\Psi(c_j),
\]
so this is a cone to $H$.  Let $t:Fx\to l$ be its supplied limiting
mediator and put $m=\Phi(t):x\to Gl$.  Then
\[
 G(p_j)m=\Phi(p_jt)=\Phi\Psi(c_j)=c_j.
\]
If $v:x\to Gl$ has the same projections, source and target naturality of
the inverse correspondence give
$p_j\Psi(v)=\Psi(G(p_j)v)=\Psi(c_j)$ for every $j$.
Limit uniqueness gives $\Psi(v)=t$, hence $v=\Phi(t)=m$.

This is an explicit formula for the new mediator: transpose the finite
input tuple, apply the given mediator, and transpose back.  All three stages
are primitive recursive.  Applying the proof in the opposite categories
proves the colimit assertion.  The same algebra applies to an infinite
indexed cone family only when the family and its universal operation have
already been presented with the requisite uniform primitive recursive
parameters; no general infinite-limit existence assertion is being made.
\end{proof}

\subsection{A constructive equivalence criterion}

A functor $F:C\to D$ is \emph{faithful} if $F(f)=F(g)$ implies $f=g$ for
parallel arrows.  Its \emph{split fullness data} are a fixed primitive
recursive operation
\[
 L_{a,a'}:\Hom_D(Fa,Fa')\longrightarrow\Hom_C(a,a')
 \quad\hbox{with}\quad F(L(h))=h.
\]
Together these say that every hom-map has a displayed inverse.  No test for
arrow equality is required.  \emph{Split essential surjectivity} means
specified primitive recursive data
\[
 s(b)\in C,\quad e_b:Fs(b)\to b,\quad d_b:b\to Fs(b),
 \qquad e_bd_b=\id_b,\quad d_be_b=\id_{Fs(b)}.
\]

\begin{theorem}[Split equivalence criterion]\label{thm:split-equivalence}
Faithfulness, split fullness, and split essential surjectivity construct a
primitive recursively presented quasi-inverse $G:D\to C$ and natural
isomorphisms $1_C\cong GF$, $FG\cong1_D$.  They may be chosen to satisfy
the triangle identities, so the result is an adjoint equivalence.
\end{theorem}
\begin{proof}
Put
\[
 G(b)=s(b),\qquad G(v:b\to b')=L_{s(b),s(b')}(d_{b'}ve_b).
\]
By split fullness,
$FG(v)=d_{b'}ve_b$.  Applying this to an identity and using faithfulness
shows $G(\id_b)=\id_{Gb}$.  For $v:b\to b'$ and $w:b'\to b''$,
\[
 F(G(w)G(v))
 =d_{b''}we_{b'}d_{b'}ve_b
 =d_{b''}wve_b=FG(wv).
\]
Faithfulness proves composition preservation.  Equality preservation of
$G$ follows by the same use of faithfulness and congruence.

The maps $e_b:FGb\to b$ are natural, since
$e_{b'}FG(v)=e_{b'}d_{b'}ve_b=ve_b$.  Their displayed inverses are $d_b$.
Define
\[
 \eta_a=L_{a,s(Fa)}(d_{Fa}),\qquad
 \theta_a=L_{s(Fa),a}(e_{Fa}).
\]
Their images under $F$ are inverse, so faithfulness proves
$\theta_a\eta_a=\id_a$ and $\eta_a\theta_a=\id_{GFa}$.  For
$u:a\to a'$, the two sides of $GF(u)\eta_a=\eta_{a'}u$ have images
\[
 d_{Fa'}F(u)e_{Fa}d_{Fa}=d_{Fa'}F(u),
\]
and hence are equal.  The inverse component families are natural as well:
multiplying the counit naturality equation by $d_{b'}$ and $d_b$ gives
$FG(v)d_b=d_{b'}v$, and multiplying the unit equation by $\theta_{a'}$
and $\theta_a$ gives $u\theta_a=\theta_{a'}GF(u)$.
Thus $\eta$ and $e$ are natural isomorphisms with displayed natural inverses.

The first triangle is
$e_{Fa}F(\eta_a)=e_{Fa}d_{Fa}=\id_{Fa}$.
For the second, apply $F$ to $G(e_b)\eta_{Gb}$ and calculate
\[
 F(G(e_b)\eta_{Gb})
 =d_be_be_{FGb}d_{FGb}=d_be_b=\id_{FGb}.
\]
Faithfulness yields $G(e_b)\eta_{Gb}=\id_{Gb}$.  All constructed operations
are finite compositions of the supplied primitive recursive operations.
\end{proof}

Conversely, for an adjoint equivalence with inverse components $\eta_a^{-1}$,
the formula $L(h)=\eta_{a'}^{-1}G(h)\eta_a$ supplies split fullness.  Indeed,
the first triangle says $F(\eta_a)=e_{Fa}^{-1}$; naturality of $e$ gives
$F(L(h))=e_{Fa'}FG(h)e_{Fa}^{-1}=h$.  If $F(f)=F(g)$, then naturality of
$\eta$ gives
$f=\eta_{a'}^{-1}GF(f)\eta_a
 =\eta_{a'}^{-1}GF(g)\eta_a=g$.
The quasi-inverse and counit supply split essential surjectivity.

The adjective ``split'' is substantive.  A mere existential preimage for
each object or arrow is not automatically a primitive recursive selector.
The usual informal phrase ``choose a preimage for every object'' is replaced
here by the displayed data that make that step legitimate.

\section{Monads and their canonical categories}\label{sec:monads}

\subsection{The monad of an adjunction}

A \emph{presented monad} consists of a fixed presented endofunctor $T:C\to C$
and natural transformations $\eta:1_C\Rightarrow T$ and
$\mu:T^2\Rightarrow T$ with primitive recursive components and laws
\[
 \mu_aT(\eta_a)=\id_{Ta}=\mu_a\eta_{Ta},\qquad
 \mu_aT(\mu_a)=\mu_a\mu_{Ta}.
\]
All transformation equalities are componentwise hom-setoid equalities.

\begin{theorem}[Every displayed adjunction induces a monad]
For $F\dashv G$ with unit $\eta$ and counit $\varepsilon$, the operations
\[
 T=GF,\qquad \mu_a=G(\varepsilon_{Fa})
\]
together with $\eta$ form a presented monad.
\end{theorem}
\begin{proof}
Functoriality of $T$ follows by composition.  Naturality of $\mu$ is the
image under $G$ of naturality of $\varepsilon$ at arrows $F(f)$.
The first unit equation is the image under $G$ of the first triangle:
\[
 \mu_aT(\eta_a)=G(\varepsilon_{Fa}F(\eta_a))=\id_{GFa}.
\]
The other is the second triangle at $Fa$:
$\mu_a\eta_{Ta}=G(\varepsilon_{Fa})\eta_{GFa}=\id_{GFa}$.
For associativity, naturality of $\varepsilon$ at
$\varepsilon_{Fa}:FGFa\to Fa$ states
\[
 \varepsilon_{Fa}FG(\varepsilon_{Fa})
 =\varepsilon_{Fa}\varepsilon_{FGFa}.
\]
Applying $G$ gives exactly $\mu_aT(\mu_a)=\mu_a\mu_{Ta}$.
\end{proof}

\subsection{Kleisli category and its adjunction}

\begin{theorem}[Kleisli construction]\label{thm:kleisli}
A presented monad $T$ determines a presented category $C_T$ with the same
object codes as $C$, arrows $a\to b$ represented by arrows $a\to Tb$ in
$C$, identity $\eta_a$, and composition
\[
 g\mathbin{\star} f=\mu_cT(g)f
 \quad(f:a\to Tb,\ g:b\to Tc).
\]
There is a displayed adjunction $F_T\dashv U_T$ that induces the original
monad.
\end{theorem}
\begin{proof}
The arrow predicate is $\Hom_C(a,Tb,f)$ and equality is the corresponding
arithmetic equality formula in $C$.  All operations have primitive recursive
code definitions.  Typing and congruence follow from those in $C$ and $T$.
The two identity equations are
\[
 \eta_b\star f=\mu_bT(\eta_b)f=f,
 \qquad
 f\star\eta_a=\mu_bT(f)\eta_a
 =\mu_b\eta_{Tb}f=f.
\]
For $h:c\to Td$, naturality of $\mu$ at $h$ gives
$T(h)\mu_c=\mu_{Td}T^2(h)$.  Thus
\begin{align*}
 h\star(g\star f)
 &=\mu_dT(h)\mu_cT(g)f\\
 &=\mu_d\mu_{Td}T^2(h)T(g)f\\
 &=\mu_dT(\mu_d)T^2(h)T(g)f\\
 &=\mu_dT(\mu_dT(h)g)f=(h\star g)\star f.
\end{align*}

Define $F_T(a)=a$, $F_T(f:a\to b)=\eta_bf$, and
$U_T(a)=Ta$, $U_T(f:a\to Tb)=\mu_bT(f)$.  Their identity laws are the
monad unit laws.  For $f:a\to b$, $g:b\to c$,
\[
 F_T(g)\star F_T(f)
 =\mu_cT(\eta_c)T(g)\eta_bf
 =T(g)\eta_bf=\eta_cgf=F_T(gf).
\]
For Kleisli arrows $f:a\to Tb$, $g:b\to Tc$, associativity and naturality
of $\mu$ give
\begin{align*}
 U_T(g\star f)
 &=\mu_cT(\mu_c)T^2(g)T(f)\\
 &=\mu_c\mu_{Tc}T^2(g)T(f)\\
 &=\mu_cT(g)\mu_bT(f)=U_T(g)U_T(f).
\end{align*}

The two sides of the hom correspondence
\[
 \Hom_{C_T}(F_Ta,b)\longleftrightarrow\Hom_C(a,U_Tb)
\]
are represented by the same codes and the same equality, so take both maps
to be the identity.  They are natural: precomposition by $F_T(u)$ in the
Kleisli category sends $f$ to
$\mu_bT(f)\eta_a u=fu$, and postcomposition by a Kleisli arrow $v$ sends
$f$ to $\mu_cT(v)f=U_T(v)f$.  The unit is the original $\eta$; the counit
at $b$ is the Kleisli arrow $Tb\to b$ represented by $\id_{Tb}$.
Its image under $U_T$ is $\mu_b$.  Moreover
$U_TF_T(f)=\mu_bT(\eta_b)T(f)=T(f)$, so the induced monad is the
original one modulo the displayed hom equalities.
\end{proof}

\subsection{Eilenberg--Moore category and its adjunction}

A $T$-algebra is a pair $(a,h)$ with $h:Ta\to a$ satisfying
\[
 h\eta_a=\id_a,\qquad h\mu_a=hT(h).
\]
A morphism $(a,h)\to(b,k)$ is an arrow $f:a\to b$ satisfying
$f h=kT(f)$.  These are arithmetical predicates on number codes, not
new comprehension axioms.

\begin{theorem}[Eilenberg--Moore construction]\label{thm:em}
The displayed $T$-algebras and their morphisms form a presented category
$C^T$.  There is a displayed adjunction $F^T\dashv U^T$, where $U^T$ forgets
the algebra operation and
\[
 F^T(a)=(Ta,\mu_a),\qquad F^T(f)=T(f).
\]
This adjunction induces the original monad.
\end{theorem}
\begin{proof}
The identity arrow is an algebra morphism because $\id_a h=hT(\id_a)$.
For algebra morphisms $f:(a,h)\to(b,k)$ and $g:(b,k)\to(c,l)$,
\[
 (gf)h=g(kT(f))=lT(g)T(f)=lT(gf).
\]
Associativity and congruence are inherited from $C$.  Thus the predicates
and operations define a presented category.

The monad laws make $(Ta,\mu_a)$ an algebra.  Naturality of $\mu$ makes
$T(f)$ an algebra morphism, so $F^T$ is a functor; its other laws and the
functor laws for $U^T$ are immediate from the corresponding supplied laws of
$T$ and $C$.

For an algebra $(b,k)$ define the correspondence
\[
 \Hom_{C^T}((Ta,\mu_a),(b,k))\longleftrightarrow\Hom_C(a,b)
\]
by $f\mapsto f\eta_a$ and $g\mapsto kT(g)$.  The second expression is
an algebra morphism because
\[
 kT(g)\mu_a=k\mu_bT^2(g)
 =kT(k)T^2(g)=kT(kT(g)).
\]
Its composite with the first map sends $g$ to
$kT(g)\eta_a=k\eta_bg=g$.
Conversely, if $f$ is an algebra morphism, the other composite is
\[
 kT(f\eta_a)=kT(f)T(\eta_a)
 =f\mu_aT(\eta_a)=f.
\]
For $u:a'\to a$, precomposition by $F^T(u)=T(u)$ changes
$f\eta_a$ to $fT(u)\eta_{a'}=f\eta_a u$.
For an algebra morphism $v:(b,k)\to(b',k')$, postcomposition changes it to
$vf\eta_a$, exactly as required.  Hence the correspondence is natural.
The unit is $\eta$ and the counit at $(b,k)$ is the algebra morphism
$k:(Tb,\mu_b)\to(b,k)$.  Its underlying arrow at the free algebra is
$\mu_a$, proving the final assertion.
\end{proof}

\subsection{Algebra limits without a preservation assumption}

\begin{theorem}[The forgetful functor creates supplied finite limits]
Let $H:J\to C^T$ be a fixed finite diagram.  If its underlying diagram in
$C$ has a supplied limit, it has a canonically supplied algebra structure
making it a limit in $C^T$.  The construction is primitive recursive and
does not require $T$ to preserve that limit.
\end{theorem}
\begin{proof}
Write $H_j=(a_j,h_j)$ and let $(p_j:l\to a_j)$ be the underlying limiting
cone.  The arrows $h_jT(p_j):Tl\to a_j$ form a cone.  Indeed, for an
algebra morphism $H(u):(a_i,h_i)\to(a_j,h_j)$,
\[
 H(u)h_iT(p_i)=h_jT(H(u))T(p_i)=h_jT(p_j).
\]
The supplied limiting operation therefore gives a unique $h:Tl\to l$
with $p_jh=h_jT(p_j)$.

To prove the unit law, compose with every projection:
\[
 p_jh\eta_l=h_jT(p_j)\eta_l
 =h_j\eta_{a_j}p_j=p_j.
\]
Limit uniqueness gives $h\eta_l=\id_l$.  To prove associativity, the two
candidate arrows $TTl\to l$ have the following projections:
\begin{align*}
 p_jh\mu_l
 &=h_jT(p_j)\mu_l
  =h_j\mu_{a_j}T^2(p_j)
  =h_jT(h_j)T^2(p_j),\\
 p_jhT(h)
 &=h_jT(p_j)T(h)
  =h_jT(p_jh)
  =h_jT(h_j)T^2(p_j).
\end{align*}
Again uniqueness proves $h\mu_l=hT(h)$.  Thus $(l,h)$ is an algebra and
its projections are algebra morphisms by their defining equations.

Finally, a cone of algebra morphisms $c_j:(x,k)\to(a_j,h_j)$ has a unique
underlying mediator $u:x\to l$.  For every $j$,
\[
 p_ju k=c_jk=h_jT(c_j)
 =h_jT(p_j)T(u)=p_jhT(u).
\]
Consequently $uk=hT(u)$ by underlying limit uniqueness, so $u$ is an
algebra morphism.  Its uniqueness in $C^T$ follows from uniqueness in $C$.
Every construction used the given finite tuple of projections, $T$, the
algebra operations, and the supplied underlying mediator; all are primitive
recursive.
\end{proof}

These results supply a substantial structural foundation: adjunctions both
produce monads and arise canonically from them, while algebraic structure
can be transported across existing universal constructions.  They do not
assert a general monadicity theorem, an adjoint functor theorem, or the
existence of arbitrary infinite colimits.  Such statements have additional
hypotheses whose arithmetic content must be separately exhibited.

\section{Pointwise Kan extensions with displayed comma colimits}\label{sec:kan}

Kan extensions organize many of the preceding universal properties. Their
existence needs care: a formula involving a colimit is not an existence
proof for that colimit. We state the exact constructive hypothesis.

Fix presented functors $K:C\to D$ and $F:C\to E$. For every object $d$
of $D$, an object of $(K\downarrow d)$ is a pair $(c,u:Kc\to d)$.
An arrow from $(c,u)$ to $(c',u')$ is $a:c\to c'$ with $u'Ka=u$.
There is a projection to $C$ and hence a diagram with value $Fc$.
Assume a colimit $L_d$ of this diagram is supplied with injections
\[
 \iota^d_{c,u}:Fc\longrightarrow L_d
\]
and its factorization and uniqueness laws. These data, and the mediator
operation on the cocones used below, must be uniform fixed PR expressions
in their numerical parameters. In particular, they respect equality of
arrow labels $u$. The cocone law is
\[
       \iota^d_{c',u'}\,Fa=\iota^d_{c,u}
                    \quad\hbox{if }u'Ka=u.                 \tag{K1}
\]
For an infinite indexing presentation, this is a schematic colimit
interface on supplied PR families; no set of all cocones is asserted.
Precisely, each fixed PR family of cocone components used in the proof,
with its displayed numerical parameters, is accompanied by a fixed PR
mediator term uniform in those parameters and by the factorization and
uniqueness proofs. The interface does not evaluate arbitrary cocone
program codes. In the finite-table case, cocones really can be supplied
as tables, and the mediator is one uniform operation on those tables.

The required equality in labels also follows from (K1): if $u\sim u'$
with the same endpoints, $\id_c$ is an arrow $(c,u)\to(c,u')$ in
the comma category, since $u'K(\id_c)\sim u$. Equation (K1),
functoriality of $F$, and the unit law therefore give
$\iota^d_{c,u}\sim\iota^d_{c,u'}$.

\begin{theorem}[Pointwise left Kan extension]\label{thm:kan}
Under these hypotheses, the assignments $d\mapsto L_d$ form a
presented functor $L:D\to E$, and there is a transformation
$\eta:F\Rightarrow LK$ such that, for every supplied presented
$G:D\to E$, composition with $\eta$ gives reciprocal constructions
\[
 \operatorname{Nat}(L,G)\ \longleftrightarrow\
 \operatorname{Nat}(F,GK).
\]
Thus $(L,\eta)$ is a left Kan extension in the presentation-wise sense.
\end{theorem}
\begin{proof}
For $v:d\to d'$, the arrows $\iota^{d'}_{c,vu}$ are a cocone over
$(K\downarrow d)$: (K1) applies because $(vu')Ka=vu$.
Let $L(v):L_d\to L_{d'}$ be its unique mediator, so
\[
                    L(v)\iota^d_{c,u}=\iota^{d'}_{c,vu}.
                                                               \tag{K2}
\]
For $v=\id$, the identity of $L_d$ has these same composites, so
uniqueness proves $L(\id)=\id$. For composable $v,w$, both $L(wv)$
and $L(w)L(v)$ have composite $\iota^{d''}_{c,wvu}$ with every
injection. Hence they are equal. The same argument shows that equal
$v$ give equal $L(v)$, so this is a functor respecting hom equality.

Define $\eta_c=\iota^{Kc}_{c,\id_{Kc}}$. If $a:c\to c'$, then
\[
 L(Ka)\eta_c=\iota^{Kc'}_{c,Ka}
            =\iota^{Kc'}_{c',\id}\,Fa=\eta_{c'}Fa,
\]
where the middle equality is (K1) for the arrow $a$ in the comma
category. Thus $\eta$ is natural.

Given $\alpha:F\Rightarrow GK$, define a cocone to $Gd$ by
\[
                   \lambda^d_{c,u}=G(u)\alpha_c.
\]
If $u'Ka=u$, its compatibility is
$G(u')\alpha_{c'}Fa=G(u')G(Ka)\alpha_c=G(u)\alpha_c$.
Let $\bar\alpha_d:L_d\to Gd$ be its mediator. To prove naturality,
compare the composites of $G(v)\bar\alpha_d$ and
$\bar\alpha_{d'}L(v)$ with $\iota^d_{c,u}$; both are
$G(vu)\alpha_c$. Uniqueness gives the naturality equation.
At the injection labelled by the identity,
\[
                  \bar\alpha_{Kc}\eta_c=\alpha_c.           \tag{K3}
\]

Conversely, suppose $\beta:L\Rightarrow G$ and
$\beta_{Kc}\eta_c=\alpha_c$. Formula (K2) gives
$\iota^d_{c,u}=L(u)\eta_c$. Therefore
\[
 \beta_d\iota^d_{c,u}
   =\beta_dL(u)\eta_c
   =G(u)\beta_{Kc}\eta_c
   =G(u)\alpha_c.
\]
Colimit uniqueness forces $\beta_d=\bar\alpha_d$ for every $d$.
This proves both inverse laws. The constructions respect transformation
equality by congruence and the same uniqueness argument.
All formulas have used fixed PR operations and finitely many logical
steps; specialization therefore yields iPRA derivations.
\end{proof}

The right Kan theorem is obtained by reversing arrows and using displayed
comma limits over $(d\downarrow K)$. If $C,D$ are finite table categories
and $F$ has finite-set values, their comma categories are finite tables.
The algorithms of Section~\ref{sec:finite} supply all required data,
so both Kan extensions exist uniformly in iPRA. For $C$ empty, the
left Kan extension has value $[0]$ and the right has value $[1]$.
The same recipes work for infinite presentations when the specific
comma (co)limits and PR mediators are supplied; their existence has not
been inferred from smallness alone.

\subsection{Universal arrows construct adjoints}
For completeness, a useful existence criterion needs much less than an
unrestricted adjoint functor theorem. Given $R:D\to C$, suppose for
each $c$ there is a supplied initial object $(Lc,\eta_c:c\to R(Lc))$
of $(c\downarrow R)$, with a uniform PR factor operation. For
$f:c\to c'$, define $Lf$ as the unique arrow satisfying
\[
                       R(Lf)\eta_c=\eta_{c'}f.
\]
The identity and composition laws follow by comparing these composites
and using initiality. The formula is also naturality of $\eta$.
For $g:c\to Rd$ the factor operation gives the unique
$\Psi(g):Lc\to d$ with $R(\Psi(g))\eta_c=g$. Its inverse is
$h\mapsto Rh\,\eta_c$. Substitution and uniqueness show naturality
in both variables, hence $L\dashv R$ by the adjunction theorem.
This is a fully constructive initial-object criterion. Merely saying
that each comma category has some initial object supplies no uniform
PR choice of the object codes or their factors; that distinction is
precisely what the hypothesis records.

\section{A constructive fixed-point theorem and an exponential obstruction}
\label{sec:lawvere}

Lawvere's diagonal theorem \cite{lawvere} provides both a further positive
result and a precise limitation on function-space claims.  Its proof is a
finite categorical calculation, so its structural content fits the present
foundation particularly well.

Suppose $C$ has a chosen terminal object and the displayed products used
below.  Fix objects $A,Y,E$, an evaluation arrow
$\operatorname{ev}:E\times A\to Y$, and $e:A\to E$.
The intended example is $E=Y^A$, but the proof needs only the specified
arrows.  For a global point $a:1\to A$, let
\[
 r_a=\operatorname{ev}\,\langle e a\,!_A,\id_A\rangle:A\to Y.
\]
Assume \emph{weak point-surjectivity}: every $f:A\to Y$ equals $r_a$ for
some $a:1\to A$.  This is an ordinary guarded arithmetic formula about
arrow codes.  In the exponential case it says that each arrow $A\to Y$
is represented by a global point in the image of $e$.

\begin{theorem}[Lawvere fixed-point theorem]\label{thm:lawvere}
Under these hypotheses, every endomorphism $t:Y\to Y$ has a global fixed
point: there exists $y:1\to Y$ with $ty=y$.  This implication is provable
in \ipra.  If weak point-surjectivity is supplied with a primitive recursive
selector $a=s(f)$, then the fixed point is given by a primitive recursive
term in $t$ and the displayed structure.
\end{theorem}
\begin{proof}
Form the diagonal arrow
\[
 h=t\,\operatorname{ev}\,\langle e,\id_A\rangle:A\to Y.
\]
Apply weak point-surjectivity once to obtain $a:1\to A$ with $h=r_a$.
Set $y=\operatorname{ev}\langle ea,a\rangle:1\to Y$.
Product uniqueness implies
$\langle e,\id_A\rangle a=\langle ea,a\rangle$, so $ha=ty$.
Terminal uniqueness gives $!_A a=\id_1$; hence
\[
 r_a a
 =\operatorname{ev}\langle ea\,!_A a,a\rangle
 =\operatorname{ev}\langle ea,a\rangle=y.
\]
The equation $h=r_a$ therefore yields $ty=y$.  Only one local existential
elimination was used, so the existential implication needs no choice
principle.  With a selector, substitute $a=s(h)$ into the displayed formula
for $y$; all resulting code operations are primitive recursive.
\end{proof}

For example, an object with two distinguished points and a pointwise
fixed-point-free negation cannot admit such a weakly point-surjective
map into its self-function object.  The assumption that negation has no
fixed \emph{global} point must be established in the particular category;
it is not a consequence of merely naming two points.

\begin{proposition}[No exponential $\N^{\N}$ among PR finite powers]
Let $\mathcal P$ be the \emph{external} category whose objects are finite
powers $\N^k$ and whose arrows are primitive recursive maps, compared by
extensional equality.  This category has finite products, but $\N^{\N}$
does not exist as an exponential object in it.  More precisely, for every
fixed PR evaluation family $V:\N^k\times\N\to\N$, a fixed PR unary function
can be constructed that differs from every member of that family.
\end{proposition}
\begin{proof}
For the fixed finite $k$, choose PR tuple decoding
$d:\N\to\N^k$ and coding $c:\N^k\to\N$ with $d(c(v))=v$.
For $k=0$ use the unique empty tuple and $c=0$.
Define
\[
 q(n)=V(d(n),n)+1.
\]
This is PR.  At $n=c(v)$ it satisfies
\[
 q(c(v))=V(v,c(v))+1\ne V(v,c(v)).
\]
Thus no tuple $v$ names $q$ in the family $V$.  The displayed inequality
is an iPRA theorem for each fixed $V,k$: it uses the coding identity and
the fact that successor has no fixed point, proved by quantifier-free
induction.

If an exponential existed, it would be one of the objects $\N^k$, and
its evaluation map would be one such PR map $V$.  The exponential property
applied to $q:\N\to\N$ would produce a global element of $\N^k$, that is,
a constant tuple $v$, with $q(n)=V(v,n)$ for every $n$.  The preceding
inequality contradicts this at the explicit index $c(v)$.
\end{proof}

The conclusion concerns this specific category and its requirement that
evaluation be PR.  It does not exclude function-space constructions with
relational evaluation, partial computation, restricted input domains,
additional clocks, or higher-type primitives.  Nor does it prohibit the
schematic Yoneda lemma: that lemma evaluates one fixed transformation term,
not a proposed PR evaluator for every PR function simultaneously.

\section{What the foundation does and does not internalize}\label{sec:boundaries}

\subsection{Why selecting structure is a mathematical hypothesis}
An existential universal property and a selected universal operation are
different interfaces. Ordinary intuitionistic logic permits the use of
one witness to prove an existential conclusion. Thus, for example, the
existence of two terminal objects implies their isomorphism without a
global choice principle. But an object-code assignment or functor requires
a function on codes. The implication
\[
  \forall x\,\exists!y\ R(x,y)
       \quad\Longrightarrow\quad
  \text{a supplied PR term }s\text{ with }R(x,s(x))
\]
is not an inference rule of iPRA. Its conclusion is also metasyntactic,
not an ordinary formula of the displayed arithmetic language. A separate
construction or a relevant extraction theorem with its precise hypotheses
would be required. Bounded search works when $R$ is PR-decidable and a PR
bound on the witness is given; our finite-table constructions provide both.

The same distinction explains the split equivalence criterion, chosen
universal elements, and the comma-colimit hypotheses for Kan extensions.
These conditions are not weakened mathematical conclusions smuggled in
as hypotheses: their exact role is to specify the computational data on
which the constructive structural theorem operates. Concrete existence
theorems must then build those data, as Section~\ref{sec:finite} does.

\subsection{A bound on isomorphism classes does not produce a skeleton}
Here is a concrete obstruction, already identified in the previous
categorical packet and developed in object-language source modules
\cite{repo-categorical,repo-two-class}. Let $H_e(t)$ be the PR bit which
reports whether machine $e$ halts in the first $t$ steps; arrange
$H_e(0)=0$. This bit is persistent: after it becomes $1$, it remains $1$.
Define the decidable relation
\[
                        s\equiv_e t\iff H_e(s)=H_e(t).
\]
Make a thin groupoid with objects $\N$, one arrow between equivalent
objects and no arrow otherwise. This is a uniform PR presentation with
at most two isomorphism classes; the assertion that three objects cannot
be pairwise inequivalent is just the two-valued pigeonhole argument.

\begin{proposition}[No uniform effective finite covering list]
There is no total computable procedure which, for every $e$, returns a
finite list meeting every isomorphism class of this groupoid.
\end{proposition}
\begin{proof}
Suppose the list is returned and let $B$ be the maximum of its entries
and $0$. If $H_e(B)=1$, the machine has halted. If $H_e(B)=0$, all listed
entries have bit $0$ by persistence. A later stage $t$ with $H_e(t)=1$
would then have no representative in the covering list, contradicting
its specification. Hence the machine never halts. The proposed procedure
would decide halting by computing the list and testing $H_e(B)$.
Such a decider is impossible: given a supposed total halting decider,
the program that halts on its own code exactly when the decider says it
does not halt yields a contradiction in either answer. This is the
ordinary external diagonal argument. No internal uniform reflection
principle is invoked.
\end{proof}

This is a computational nonexistence theorem, not a claim that no
particular groupoid has a skeleton. A supplied finite covering list
changes the situation. With decidable equivalence, choose the first
cover entry equivalent to the input and remove later equivalent entries
by a bounded scan. The representative is found because of the covering
hypothesis, and is constant on equivalence classes by transitivity and
the least-index definition. For these thin groupoids, the unique arrows
give the comparison isomorphisms. In a general category, one must also
supply or compute the isomorphism arrows, as a finite table category
allows. ProveIt's \texttt{FiniteCoverSkeleton} records this positive
construction \cite{repo-skeleton}.

\subsection{Large collections, function spaces and infinitary theorems}
The phrase ``category of all sets'' cannot be read literally as a category
of natural-number codes containing every ordinary set. This report
internalizes specified arithmetic presentations, with semantic
interpretations external when needed. It does not claim that arbitrary
sets, arbitrary real numbers, or every classical space has a code.
It nevertheless allows infinite categories and undecidable hom equality.

Likewise, the schematic Yoneda embedding does not by itself construct a
single internal object containing every presheaf or every natural
transformation. A particular uniformly evaluable coded family can be
used, and a finite functor category can be constructed completely.
General relational representations remain possible, but require their
own totality and application interface. The distinction is supported
by the explicit diagonal obstruction of Section~\ref{sec:lawvere} and
the existing evaluator module \cite{repo-evaluator}.

The main results proved here cover the usual structural core of a first
course, extended by explicit density, monads, Kan extensions, idempotents
and diagonal arguments. They do not assert all theorems in every branch
of category theory. In particular:
\begin{itemize}
\item General adjoint functor theorems need their solution-set and
      (co)completeness assumptions, together with constructive witness
      information; the initial-comma-object criterion proved here
      isolates the exact positive interface.
\item Arbitrary small limits and colimits do not follow from the finite
      constructions. Even the category of finite sets cannot contain
      all infinite products of nontrivial finite sets.
\item Beck's general monadicity theorem, general monoidal coherence,
      topos representation theorems, arbitrary localizations and
      higher-category theorems are not claimed proved by this packet.
      The relevant notions need additional definitions and hypotheses,
      not an invocation of the word ``constructive''.
\item No universal iPRA theorem saying that every accepted iPRA proof
      is true has been added. Proof transformation operates on concrete
      derivations; it is distinct from internal soundness or consistency
      reflection.
\end{itemize}

The obstruction is therefore not categorical abstraction itself. The
finite structural proofs work with arithmetic equality of any complexity.
The additional strength enters when a proof collects unrestricted objects,
chooses witnesses without a displayed operation, evaluates all PR programs
uniformly, or inducts on a predicate outside the allowed scheme.

\section{Lean architecture, repository reuse and completion path}\label{sec:lean}

\subsection{The target is a derivation in iPRA}
The relevant output type is
\[
            \Gamma\vdash[\mathsf{PRAAx}]\varphi,
       \qquad \varphi:\mathsf{Fm}\;\mathsf{PTm}\;n,
\]
not merely a proposition proved in Lean's ambient logic. The supplied
source files use the actual ProveIt notation for this judgment.
Here \texttt{PTm} is scoped term syntax, \texttt{Fm} is scoped
first-order formula syntax, and \texttt{Prf} is an inductive
intuitionistic natural-deduction derivation. \texttt{PRAAx} contains
exactly the arithmetic axioms of Section~\ref{sec:foundation}.

The existing syntax is arity-indexed. In schematic Lean notation:
\begin{verbatim}
inductive PR : Nat -> Type
  | zero : PR 0
  | succ : PR 1
  | proj {k} (i : Fin k) : PR k
  | comp {m k} (f : PR m) (gs : Fin m -> PR k) : PR k
  | prec {k} (f : PR k) (g : PR (k + 2)) : PR (k + 1)

inductive PTm (n : Nat) : Type
  | var (i : Fin n) : PTm n
  | app {k} (f : PR k) (ts : Fin k -> PTm n) : PTm n
\end{verbatim}
The finite functions in these constructors are tuples of syntax in the
metatheory. They are not arbitrary functions admitted into iPRA.
Formula constructors are equality of two terms, false, conjunction,
disjunction, implication, and the two quantifiers; the latter reduce
the number of free variables by one. Substitution shifts indices under
binders. \texttt{Fm.QF} recursively excludes both quantifier constructors.
The induction constructor requires an explicit proof of this property.
The appendix source closure includes these definitions and their
substitution lemmas, rather than replacing them with unspecified axioms.

\subsection{Essential data records}
The definition \texttt{PresentedCategory.Pres} uses the following exact
field types. Its associated \texttt{Laws} record contains derivation
constructors, not assertions that hold only in a semantic interpretation.
\begin{center}
\begin{tabular}{lll}
\toprule
Field & Syntax type & Intended arguments\\
\midrule
\texttt{Obj} & \texttt{Fm PTm 1} & $a$\\
\texttt{Hom} & \texttt{Fm PTm 3} & $a,b,f$\\
\texttt{Eqv} & \texttt{Fm PTm 4} & $a,b,f,g$\\
\texttt{idt} & \texttt{PTm 1} & $a$\\
\texttt{cmp} & \texttt{PTm 5} & $a,b,c,g,f$\\
\bottomrule
\end{tabular}
\end{center}
For example, its associativity field maps derivations typing $f,g,h$
to a derivation of the appropriate instantiated \texttt{Eqv} formula.
The human proof's endpoint guards are retained explicitly when needed;
the upstream minimal record does not supply every possible endpoint
closure rule automatically.

\texttt{Adjunction.Fun} supplies object and arrow terms of arities one
and three. A general natural transformation has a component term of
arity one and object-language typing and naturality laws. A family with
additional numerical parameters adds those parameters to the term arity.
\texttt{Yoneda.Fam} has the formulas \texttt{Par}, \texttt{Car},
\texttt{Sim} of arities one, three and four, and \texttt{act : PTm 5}.
Its law record supplies the presheaf action and congruence laws used in
the Yoneda calculation. For subsequent constructions requiring all
setoid equivalence or carrier-closure laws, those laws must also be
supplied; a minimal upstream record is not silently strengthened.

For a fixed diagram with objects $D_0,\ldots,D_{r-1}$ and arrows $D(u)$,
the cone formula at vertex $x$ and a coded component tuple $c$ is the
finite conjunction
\[
 \operatorname{Cone}_D(x,c)\equiv
 \bigwedge_i H_C(x,D_i,c_i)\ \wedge\
 \bigwedge_{u:i\to j}E_C(x,D_j,D(u)c_i,c_j).
\]
A selected limit record contains a term $L$, projection terms $p_i$,
a mediator term $m(x,c)$, and derivations of: the cone laws; typing of
$m$ under the cone premise; $p_i m=c_i$; and
\[
 H_C(x,L,v)\wedge\bigwedge_i E_C(x,D_i,p_i v,c_i)
           \ \longrightarrow\ E_C(x,L,v,m(x,c)).
\]
Each fixed conjunction is ordinary formula syntax. Products, equalizers
and colimits are specializations or duals. Varying diagram length is
handled by the finite-table Boolean representation, not by this fixed
conjunction as though it were a PR truth evaluator for arbitrary formulas.

A monad record adds fixed terms for $T$, $\eta$ and $\mu$, typing and
naturality derivations, and the three monad equations. The Kleisli
presentation has hom formula $H_C(a,Tb,f)$ and composition term
$\mu_cT(g)f$. The Eilenberg--Moore object formula on $(a,h)$ is
\[
 O_C(a)\wedge H_C(Ta,a,h)\wedge h\eta_a=\id_a
                         \wedge h\mu_a=hT(h),
\]
with its displayed hom predicate from Section~\ref{sec:monads}.
The Karoubi object formula uses $e^2=e$, and its identity term returns
$e$. These arithmetical formulas may be undecidable, but their identity
and composition constructors are PR. To complete their Lean proofs,
one expands the finite calculations in the report into the already
available equality and substitution derivation constructors.

\subsection{What was actually inspected in ProveIt}
All source reuse is pinned to commit
\texttt{72d37f6d06fa7b5bd0a46732060446d1cffa3f2d}, retrieved on
20 November 2026 through the GitHub connector. The repository specifies
Lean \texttt{v4.32.0} and Mathlib \texttt{v4.32.0}.
The package contains 104 upstream Lean files forming the complete
in-repository import closure of the selected categorical modules, their
license, and a source manifest. Mathlib and the Lean executable are
external build dependencies.

\begin{longtable}{@{}>{\raggedright\arraybackslash}p{0.30\linewidth}>{\raggedright\arraybackslash}p{0.63\linewidth}@{}}
\toprule
Module & Inspected object-language content and limitations\\
\midrule
\endhead
\texttt{PresentedCategory} & Category formulas and PR terms; inverse
uniqueness, sections and retractions, terminal uniqueness; a concrete
additive-monoid instance.\\
\texttt{Yoneda} & Presheaf and transformation law records;
\texttt{yoneda}, \texttt{yoneda\_closed}, \texttt{repr\_natural},
\texttt{repr\_eval}, \texttt{rep\_laws}, representable specialization.\\
\texttt{Adjunction} & Functor and adjunction records; both hom-map
round trips; construction of the unit and counit from a natural
correspondence, with naturality and triangles.\\
\texttt{SplitEquivalence} & Explicit hom lifting, object representatives,
quasi-inverse, inverse equations and naturality.\\
\texttt{FiniteLimits} & Rank/select construction at a fixed diagram
shape; assumes a retention test with soundness and completeness and
positive size conditions. The general table-built test is not supplied.\\
\texttt{FiniteSetClosed} & Evaluation, currying and their round trips
under positive-cardinality hypotheses. Empty cases need extension.\\
\texttt{FiniteCoequalizer} & Sequential label merging and its
factorization calculations in the source's stated finite-set encoding.\\
\texttt{KanFormula} & Representative-level left and right Kan formulas,
their compatibility and naturality; not an unrestricted existence
theorem for arbitrary comma (co)limits.\\
\texttt{SetoidInterpretation} & Domain- and equality-respecting
translation of object-language derivations.\\
\texttt{UniversalEvaluator}\linebreak\texttt{Obstruction} & An external no-universal-PR-
evaluator theorem and, separately, an actual iPRA diagonal sentence.\\
\texttt{FiniteCoverSkeleton}, \texttt{TwoClassObstruction} & Positive
bounded representative selection and the precise obstruction to
extracting it merely from a bound on equivalence classes.\\
\bottomrule
\end{longtable}

The repository describes its \texttt{Derivations} library as kernel-checked.
This packet has inspected those pinned sources but has \emph{not} rebuilt
them locally. A proof-shaped source declaration is not by itself a fresh
build receipt. Nor should an older research packet's description of
``uncompiled sketches'' be applied to a newer \texttt{Derivations} module
without reading the latter. These distinctions matter when reporting reuse
\cite{repo-readme,repo-presented,repo-yoneda,repo-adjunction}.

\subsection{New Lean adapters and a useful interface refinement}
The source \texttt{Yoneda.NatLaws} asks one term
$\alpha(p,A,x,f)$ to obey its laws uniformly for every family parameter
$p$ and representing object $A$. This is a valid sufficient interface,
but stronger than the data of one transformation $h_A\Rightarrow P_p$
at fixed $p,A$. It should not be mistaken for that local interface.

The new \texttt{CategoryIPRA/LocalYoneda.lean} defines
\texttt{AtLaws} at fixed $p,A$ and object-language context $\Gamma$.
It repeats the short Yoneda derivation with these local premises and
provides conversion from the upstream uniform record.
\texttt{Representability.lean} derives faithfulness and the two inverse
equations from a specified natural isomorphism, using this localized
record. \texttt{Density.lean} derives value typing, canonical-value
evaluation, generator round trip and reindexing for the density
construction. Their conclusions are actual iPRA derivation judgments;
they introduce no universal evaluator and no new axiom declaration.

\textbf{Verification status.} All three new Lean modules are uncompiled
integration drafts. No Lean executable was available in the generation
environment. They contain no \texttt{sorry} or \texttt{admit}, but their
elaboration and kernel checking remain to be performed. The completed
human proofs are separate evidence; no statement in this report treats
the drafts as a completed mechanization. The source-structure audit
parsed the Lake configuration and found no missing in-repository modules
among 192 imports across 110 Lean source files. External import names lie
under the specified Mathlib dependency; its installed module files were
not checked because the build environment was unavailable.

\subsection{A concrete path to completion}
The archive's \texttt{lean/formalization\_roadmap.md} gives record fields,
dependencies and coding obligations in more detail. The following order
keeps each step mathematically reviewable.
\begin{enumerate}
\item Build the pinned source closure and the three new adapters; resolve
elaboration issues and record the actual kernel and axiom-audit output.
Close the desired local derivations by the existing universal-closure
machinery, constructing law records under the required contexts.
\item Add generic natural transformations, selected cones and universal
operations as syntax records. Port their finite equational proofs using
the supplied category and functor law constructors.
\item Complete density's tuple coding and cocone record. The difficult
logical point is already removed by direct equality of values; no
unbounded generated-equivalence closure is needed.
\item Port the base-$(m+1)$ finite tables, the Boolean diagram tester,
rank/select, and all zero-cardinality cases. State every loop invariant
as a concrete quantifier-free formula. The executable model provides
reference outputs, not proofs of these invariants.
\item Encode the monad, Kleisli, algebra and idempotent formulas and
operations exactly as above. Each law is a finite equation chain;
algebra limits use the selected limiting mediator and its uniqueness.
\item Encode the parameterized comma diagrams and their selected
(co)limits. Define the Kan functor and mediator by the formulas in
Section~\ref{sec:kan}. Preserve the explicit existence assumptions.
\item Audit every theorem as either a closed iPRA derivation, a derivation
under named hypotheses, or an external schema. In particular, do not
replace code-level PR certificates by unrestricted semantic Lean
functions, and do not infer an internal soundness principle from a
metatheoretic proof transformation.
\end{enumerate}

With the pinned toolchain and Mathlib available, the intended commands are
\begin{verbatim}
cd lean
lake update
lake exe cache get
lake build CategoryIPRA
lake env lean CategoryIPRA/Audit.lean
\end{verbatim}
These are reproduction instructions. They have not been presented as
successful executions in this environment.

\subsection{Executable validation and its limits}
The exact-arithmetic model in \texttt{code/finite\_models.py} uses integer
tables, including empty tables and base-$(m+1)$ encodings. An actual run
passed 22,016 exponential beta cases and 22,016 eta cases; 1,544 product
beta cases and 1,544 eta cases; 60 valid and 84 invalid table-code cases;
and eight limit/colimit example diagrams, whose universal properties were
exhaustively checked against test objects of cardinality at most three.
The diagrams include the empty diagram, empty objects, an empty equalizer,
a pullback and a quotient requiring transitive identifications.

The powerset presheaf on the finite full subcategory with cardinalities
$0,1,2$ gives a more discriminating Yoneda test: all 1,024 possible
component families were checked, and precisely four were natural,
corresponding to the four elements of the powerset of $[2]$.
Two further examples use the group and zero monoid of order two.
The option monad's Kleisli operations passed 23 unit cases and 2,457
associativity cases.

The executable colimit implementation uses bounded Warshall closure;
the human proof gives sequential merging. Both realize the same
finite-path quotient. This independence is useful for catching example
errors, but tests of finitely many models do not prove an unrestricted
theorem or verify an iPRA proof checker. The JSON receipt records the
actual finite scope and all counts.

\clearpage
\phantomsection
\addcontentsline{toc}{section}{References}
\begin{thebibliography}{99}
\bibitem{riehl} Emily Riehl. \emph{Category Theory in Context}. Dover, 2016; author-hosted updated edition, consulted 20 November 2026. Standard background and scope; the arithmetic presentations and proofs in this report are developed explicitly. \url{https://emilyriehl.github.io/files/context.pdf}.
\bibitem{awodey} Steve Awodey. \emph{Categorical Logic}, course notes, Appendix A: Category Theory, draft dated 16 January 2024. \url{https://awodey.github.io/catlog/notes/catlog0.pdf}.
\bibitem{lawvere} F. William Lawvere. Diagonal arguments and cartesian closed categories. \emph{Lecture Notes in Mathematics} 92 (1969), 134--145. Reprinted with author commentary, \emph{Reprints in Theory and Applications of Categories} 15 (2006), 1--13. \url{https://www.tac.mta.ca/tac/reprints/articles/15/tr15abs.html}.
\bibitem{wiki} Wikipedia contributors. Category theory. User-supplied overview, consulted 20 November 2026; used to identify scope, not as a substitute for proofs. \url{https://en.wikipedia.org/wiki/Category_theory}.
\bibitem{repo-readme} ProveIt contributors. Intuitionistic primitive recursive arithmetic: repository overview. Pinned commit \texttt{72d37f6d06fa}. \url{https://github.com/VladimirReshetnikov/ProveIt/blob/72d37f6d06fa7b5bd0a46732060446d1cffa3f2d/Logic/PrimitiveRecursiveArithmetic/README.md}.
\bibitem{repo-pr} ProveIt contributors. PR syntax, terms, and axioms of iPRA. Pinned commit \texttt{72d37f6d06fa}. \url{https://github.com/VladimirReshetnikov/ProveIt/blob/72d37f6d06fa7b5bd0a46732060446d1cffa3f2d/Logic/PrimitiveRecursiveArithmetic/ConciseIPRA/Lean/ConciseIPRA/PR.lean}.
\bibitem{repo-presented} ProveIt contributors. Presented categories and elementary iPRA derivations. Pinned commit \texttt{72d37f6d06fa}. \url{https://github.com/VladimirReshetnikov/ProveIt/blob/72d37f6d06fa7b5bd0a46732060446d1cffa3f2d/Logic/PrimitiveRecursiveArithmetic/Derivations/Lean/IPRADerivations/Foundations/PresentedCategory.lean}.
\bibitem{repo-yoneda} ProveIt contributors. Yoneda derivations. Pinned commit \texttt{72d37f6d06fa}. \url{https://github.com/VladimirReshetnikov/ProveIt/blob/72d37f6d06fa7b5bd0a46732060446d1cffa3f2d/Logic/PrimitiveRecursiveArithmetic/Derivations/Lean/IPRADerivations/Foundations/Yoneda.lean}.
\bibitem{repo-adjunction} ProveIt contributors. Adjunction derivations. Pinned commit \texttt{72d37f6d06fa}. \url{https://github.com/VladimirReshetnikov/ProveIt/blob/72d37f6d06fa7b5bd0a46732060446d1cffa3f2d/Logic/PrimitiveRecursiveArithmetic/Derivations/Lean/IPRADerivations/Foundations/Adjunction.lean}.
\bibitem{repo-setoid} ProveIt contributors. Setoid interpretation of derivations. Pinned commit \texttt{72d37f6d06fa}. \url{https://github.com/VladimirReshetnikov/ProveIt/blob/72d37f6d06fa7b5bd0a46732060446d1cffa3f2d/Logic/PrimitiveRecursiveArithmetic/Derivations/Lean/IPRADerivations/Foundations/SetoidInterpretation.lean}.
\bibitem{repo-evaluator} ProveIt contributors. Universal evaluator obstruction. Pinned commit \texttt{72d37f6d06fa}. \url{https://github.com/VladimirReshetnikov/ProveIt/blob/72d37f6d06fa7b5bd0a46732060446d1cffa3f2d/Logic/PrimitiveRecursiveArithmetic/Derivations/Lean/IPRADerivations/Foundations/UniversalEvaluatorObstruction.lean}.
\bibitem{repo-two-class} ProveIt contributors. Two-class obstruction. Pinned commit \texttt{72d37f6d06fa}. \url{https://github.com/VladimirReshetnikov/ProveIt/blob/72d37f6d06fa7b5bd0a46732060446d1cffa3f2d/Logic/PrimitiveRecursiveArithmetic/Derivations/Lean/IPRADerivations/Foundations/TwoClassObstruction.lean}.
\bibitem{repo-skeleton} ProveIt contributors. Finite-cover skeleton construction. Pinned commit \texttt{72d37f6d06fa}. \url{https://github.com/VladimirReshetnikov/ProveIt/blob/72d37f6d06fa7b5bd0a46732060446d1cffa3f2d/Logic/PrimitiveRecursiveArithmetic/Derivations/Lean/IPRADerivations/Foundations/FiniteCoverSkeleton.lean}.
\bibitem{repo-categorical} ProveIt contributors. Earlier categorical methods research packet, including category scope and obstruction sections. Pinned commit \texttt{72d37f6d06fa}. \url{https://github.com/VladimirReshetnikov/ProveIt/blob/72d37f6d06fa7b5bd0a46732060446d1cffa3f2d/Logic/PrimitiveRecursiveArithmetic/docs/proofs/foundations/ipra_categorical_methods/README.md}.
\end{thebibliography}


\end{document}

